\chaptertitle{Little Harmonic Labyrinth}

                     Little Harmonic Labyrinth

  The Tortoise and Achilles are spending a day at Coney Island After buying a
  couple of cotton candies, they decide to take a ride on the Ferris wheel.

Tortoise: This is my favorite ride. One seems to move so far, and
reality one gets nowhere.
Achilles: I can see why it would appeal to you. Are you all strapped in?
Tortoise: Yes, I think I've got this buckle done. Well, here we go.
Achilles: You certainly are exuberant today.
Tortoise: I have good reason to be. My aunt, who is a fortune-teller me that a stroke of
        Good Fortune would befall me today. So I am tingling with anticipation.
Achilles: Don't tell me you believe in fortune-telling!
Tortoise: No ... but they say it works even if you don't believe ii
Achilles: Well, that's fortunate indeed.
Tortoise: Ah, what a view of the beach, the crowd, the ocean, the city. . .
Achilles: Yes, it certainly is splendid. Say, look at that helicopter there. It seems to be
        flying our way. In fact it's almost directly above us now.
Tortoise: Strange-there's a cable dangling down from it, which is very close to us. It's
        coming so close we could practically grab it
Achilles: Look! At the end of the line there's a giant hook, with a note

(He reaches out and snatches the note. They pass by and are on their z down.)

Tortoise: Can you make out what the note says?
Achilles: Yes-it reads, "Howdy, friends. Grab a hold of the hook time around, for an
       Unexpected Surprise."
Tortoise: The note's a little corny but who knows where it might lead, Perhaps it's got
       something to do with that bit of Good Fortune due me. By all means, let's try it!
Achilles: Let's!

  (On the trip up they unbuckle their buckles, and at the crest of the ride, grab for the
  giant hook. All of a sudden they are whooshed up by the ca which quickly reels
  them skyward into the hovering helicopter. A It strong hand helps them in.)

Voice: Welcome aboard-Suckers.
Achilles: Wh-who are you?




Little Harmonic Labyrinth                                                               111

Voce: Allow me to introduce myself. I am Hexachlorophene J. Goodforttune, Kidnapper
       At-Large, and Devourer of Tortoises par Excellence, at your service.
Tortoise: Gulp!
Achilles (whispering to his friend): Uh-oh-I think that this "Goodfortune" is not exactly
       what we'd anticipated. (To Goodfortune) Ah-if I may be so bold-where are you
       spiriting us off to?
Goodfortune: Ho ho! To my all-electric kitchen-in-the-sky, where I will prepare THIS
       tasty morsel-(leering at the Tortoise as he says this)-in a delicious pie-in-the-sky!
       And make no mistake-it's all just for my gobbling pleasure! Ho ho ho!
Achilles: All I can say is you've got a pretty fiendish laugh.
Goodfortune (laughing fiendishly): Ho ho ho! For that remark, my friend, you will pay
       dearly. Ho ho!
Achilles: Good grief-I wonder what he means by that!
Goodfortune: Very simple-I've got a Sinister Fate in store for both of you! Just you wait!
       Ho ho ho! Ho ho ho!
Achilles: Yikes!
Goodfortune: Well, we have arrived. Disembark, my friends, into my fabulous all-electric
       kitchen-in-the-sky.

(They walk inside.)

    Let me show you around, before I prepare your fates. Here is my bedroom. Here is
    my study. Please wait here for me for a moment. I've got to go sharpen my knives.
    While you're waiting, help yourselves to some popcorn. Ho ho ho! Tortoise pie!
    Tortoise pie! My favorite kind of pie! (Exit.)

Achilles: Oh, boy-popcorn! I'm going to munch my head off!
Tortoise: Achilles! You just stuffed yourself with cotton candy! Besides, how can you
       think about food at a time like this?
Achilles: Good gravy-oh, pardon me-I shouldn't use that turn of phrase, should I? I mean
       in these dire circumstances ... Tortoise: I'm afraid our goose is cooked.
Achilles: Say-take a gander at all these books old Goodfortune has in his study. Quite a
       collection of esoterica: Birdbrains I Have Known; Chess and Umbrella-Twirling
       Made Easy; Concerto for Tapdancer and Orchestra ... Hmmm.
Tortoise: What's that small volume lying open over there on the desk, next to the
       dodecahedron and the open drawing pad?
Achilles: This one? Why, its title is Provocative Adventures of Achilles and the Tortoise
       Taking Place in Sundry Spots of the Globe. Tortoise: A moderately provocative
       title.
Achilles: Indeed-and the adventure it's opened to looks provocative. It's called "Djinn and
       Tonic".
Tortoise: Hmm ... I wonder why. Shall we try reading it? I could take the Tortoise's part,
       and you could take that of Achilles.




Little Harmonic Labyrinth                                                               112

Achilles: I’m game. Here goes nothing . . .

       (They begin reading "Djinn and Tonic".)

                 (Achilles has invited the Tortoise over to see his collection of prints by
                 his favorite artist, M. C. Escher.)

           Tortoise: These are wonderful prints, Achilles.
           Achilles: I knew you would enjoy seeing them. Do you have any particular
                  favorite?
           Tortoise: One of my favorites is Convex and Concave, where two internally
                  consistent worlds, when juxtaposed, make a completely inconsistent
                  composite world. Inconsistent worlds are always fun places to visit,
                  but I wouldn't want to live there.
           Achilles: What do you mean, "fun to visit"? Inconsistent worlds don't EXIST,
                  so how can you visit one?
           Tortoise: I beg your pardon, but weren't we just agreeing that in

           this Escher picture, an inconsistent world is portrayed?
           Achilles: Yes, but that's just a two-dimensional world-a fictitious world-a
                   picture. You can't visit that world.
           Tortoise: I have my ways ...
           Achilles: How could you propel yourself into a flat picture-universe?
           Tortoise: By drinking a little glass of PUSHING-POTION. That does the
                   trick.
           Achilles: What on earth is pushing-potion?
           Tortoise: It's a liquid that comes in small ceramic phials, and which, when
                   drunk by someone looking at a picture, "pushes'' him right into the
                   world of that picture. People who aren't aware of the powers of
                   pushing-potion often are pretty surprised by the situations they wind
                   up in.
           Achilles: Is there no antidote? Once pushed, is one irretrievably lost?
           Tortoise: In certain cases, that's not so bad a fate. But there is, in fact, another
                   potion-well, not a potion, actually, but an elixir-no, not an elixir, but
                   a-a
           Tortoise: He probably means "tonic".
           Achilles: Tonic?
           Tortoise: That's the word I was looking for! "POPPING-TONIC" iu what it's
                   called, and if you remember to carry a bottle of it in your right hand as
                   you swallow the pushing-potion, it too will be pushed into the picture;
                   then, whenever you get a hanker ing to "pop" back out into real life,
                   you need only take a swallow of popping-tonic, and presto! You're
                   back in the rea. world, exactly where you were before you pushed
                   yourself in.
           Achilles: That sounds very interesting. What would happen it you took some
                   popping-tonic without having previously pushed yourself into a
                   picture?


Little Harmonic Labyrinth                                                                 113

          Tortoise: I don’t precisely know, Achilles, but I would be rather wary of
                 horsing around with these strange pushing and popping liquids. Once I
                 had a friend, a Weasel, who did precisely what you suggested-and no
                 one has heard from him since.
          Achilles: That's unfortunate. Can you also carry along the bottle of pushing-
                 potion with you?
          Tortoise: Oh, certainly. Just hold it in your left hand, and it too will get
                 pushed right along with you into the picture you're looking at.
          Achilles: What happens if you then find a picture inside the picture which you
                 have already entered, and take another swig of pushing-potion?
          Tortoise: Just what you would expect: you wind up inside that picture-in-a-
                 picture.
          Achilles: I suppose that you have to pop twice, then, in order to extricate
                 yourself from the nested pictures, and re-emerge back in real life.
          Tortoise: That's right. You have to pop once for each push, since a push takes
                 you down inside a picture, and a pop undoes that.
          Achilles: You know, this all sounds pretty fishy to me . . . Are you sure you're
                 not just testing the limits of my gullibility?
          Tortoise: I swear! Look-here are two phials, right here in my pocket.
                 (Reaches into his lapel pocket, and pulls out two rather large
                 unlabeled phials, in one of which one can hear a red liquid sloshing
                 around, and in the other of which one can hear a blue liquid sloshing
                 around.) If you're willing, we can try them. What do you say?
          Achilles: Well, I guess, ahm, maybe, ahm ...
          Tortoise: Good! I knew you'd want to try it out. Shall we push ourselves into
                 the world of Escher's Convex and Concave?
          Achilles: Well, ah, .. .
          Tortoise: Then it's decided. Now we've got to remember to take along this
                 flask of tonic, so that we can pop back out. Do you want to take that
                 heavy responsibility, Achilles?
          Achilles: If it's all the same to you, I'm a little nervous, and I'd prefer letting
                 you, with your experience, manage the operation.
          Tortoise: Very well, then.

             (So saying, the Tortoise pours two small portions of pushing-potion. Then
             he picks up the flask of tonic and grasps it firmly in his right hand, and
             both he and Achilles lift their glasses to their lips.)

          Tortoise: Bottoms up!

          (They swallow.)




Little Harmonic Labyrinth                                                               114

              FIGURE 23. Convex and Concave, by M. C. Escher (lithograph, 1955).


          Achilles: That's an exceedingly strange taste.
          Tortoise: One gets used to it.
          Achilles: Does taking the tonic feel this strange? Tortoise: Oh, that's quite
                 another sensation. Whenever you taste the tonic, you feel a deep sense
                 of satisfaction, as if you'd been waiting to taste it all your life.
                 Achilles: Oh, I'm looking forward to that. Tortoise: Well, Achilles,
                 where are we?
          Achilles (taking cognizance of his surroundings): We're in a little gondola,
                 gliding down a canal! I want to get out. Mr.Gondolier, please let us
                 out here.

                  (The gondolier pays no attention to this request.)

          Tortoise: He doesn't speak English. If we want to get out here, we'd better just
                 clamber out quickly before he




Little Harmonic Labyrinth                                                             115

          Enters the sinister “Tunnel of Love”; just ahead of us.

             (Achilles, his face a little pale scrambles out in a split second and then
             pulls his slower friend out.)

          Achilles: I didn't like the sound of that place, somehow. I'm glad we got out
                 here. Say, how do you know so much about this place, anyway? Have
                 you been here before?
          Tortoise: Many times, although I always came in from other Escher pictures.
                 They're all connected behind the frames, you know. Once you're in
                 one, you can get to any other one.
          Achilles: Amazing! Were I not here, seeing these things with my own eyes,
                 I'm not sure I'd believe you. (They wander out through a little arch.)
                 Oh, look at those two cute lizards!
          Tortoise: Cute? They aren't cute-it makes me shudder just to think of them!
                 They are the vicious guardians of that magic copper lamp hanging
                 from the ceiling over there. A mere touch of their tongues, and any
                 mortal turns to a pickle.
          Achilles: Dill, or sweet?
          Tortoise: Dill.
          Achilles: Oh, what a sour fate! But if the lamp has magical powers, I would
                 like to try for it.
          Tortoise: It's a foolhardy venture, my friend. I wouldn't risk it.
          Achilles: I'm going to try just once.

             (He stealthily approaches the lamp, making sure not to awaken the
             sleeping lad nearby. But suddenly, he slips on a strange shell-like
             indentation in the floor, and lunges out into space. Lurching crazily, he
             reaches for anything, and manages somehow to grab onto the lamp with
             one hand. Swinging wildly, with both lizards hissing and thrusting their
             tongues violently out at him, he is left dangling helplessly out in the middle
             of space.)

          Achilles: He-e-e-elp!

             (His cry attracts the attention of a woman who rushes downstairs and
             awakens the sleeping boy. He takes stock of the situation, and, with a
             kindly smile on his face, gestures to Achilles that all will be well. He shouts
             something in a strange guttural tongue to a pair of trumpeters high up in
             windows, and immediately,




Little Harmonic Labyrinth                                                               116

             Weird tones begin ringing out and making beats each other. The sleepy
             young lad points at the lizards, and Achilles sees that the music is having a
             strong soporific effect on them. Soon, they are completely unconscious.
             Then the helpful lad shouts to two companions climbing up ladders. They
             both pull their ladders up and then extend them out into space just
             underneath the stranded Achilles, forming a sort of bridge. Their gestures
             make it clear that Achilles should hurry and climb on. But before he does
             so, Achilles carefully unlinks the top link of the chain holding the lamp, and
             detaches the lamp. Then he climbs onto the ladder-bridge and the three
             young lads pull him in to safety. Achilles throws his arms around them and
             hugs them gratefully.)

          Achilles: Oh, Mr. T, how can I repay them?
          Tortoise: I happen to know that these valiant lads just love coffee, and down
                 in the town below, there's a place where they make an incomparable
                 cup of espresso. Invite them for a cup of espresso! Achilles: That
                 would hit the spot.

             (And so, by a rather comical series of gestures, smiles, and words, Achilles
             manages to convey his invitation to the young lads, and the party of five
             walks out and down a steep staircase descending into the town. They reach
             a charming small cafe, sit down outside, and order five espressos. As they
             sip their drinks, Achilles remembers he has the lamp with him.)

          Achilles: I forgot, Mr. Tortoise-I've got this ma; lamp with me! But-what's
                 magic about it? Tortoise: Oh, you know, just the usual-a genie.
          Achilles: What? You mean a genie comes out when you rub it, and grants you
                 wishes?
          Tortoise: Right. What did you expect? Pennies fry heaven?
          Achilles: Well, this is fantastic! I can have any wish want, eh? I've always
                 wished this would happen to me ...

             (And so Achilles gently rubs the large letter `L' which is etched on the
             lamp's copper surface ... Suddenly a huge puff of smoke appears, and in the
             forms of the smoke the five friends can make out a weird, ghostly figure
             towering above them.)




Little Harmonic Labyrinth                                                              117

            I
          Genie: Hello, my friends – and thanks ever so much for rescuing my Lamp
                 from the evil Lizard-Duo.

             (And so saying, the Genie picks up the Lamp, and stuffs it into a pocket
             concealed among the folds of his long ghostly robe which swirls out of the
             Lamp.)

          As a sign of gratitude for your heroic deed, I would like to offer you, on the
                 part of my Lamp, the opportunity to have any three of your wishes
                 realized.
          Achilles: How stupefying! Don't you think so, Mr. T?
          Tortoise: I surely do. Go ahead, Achilles, take the first wish.
          Achilles: Wow! But what should I wish? Oh, I know! It's what I thought of
                 the first time I read the Arabian Nights (that collection of silly (and
                 nested) tales)-I wish that I had a HUNDRED wishes, instead of just
                 three! Pretty clever, eh, Mr. T? I bet YOU never would have thought
                 of that trick. I always wondered why those dopey people in the stories
                 never tried it themselves.
          Tortoise: Maybe now you'll find out the answer.
          Genie: I am sorry, Achilles, but I don't grant metawishes.
          Achilles: I wish you'd tell me what a "meta-wish" is!
          Genie: But THAT is a meta-meta-wish, Achilles-and I don't grant them,
                 either. Achilles: Whaaat? I don't follow you at all.
          Tortoise: Why don't you rephrase your last request, Achilles?
          Achilles: What do you mean? Why should I?
          Tortoise: Well, you began by saying "I wish". Since you're just asking for
                 information, why don't you just ask a question?
          Achilles: All right, though I don't see why. Tell me, Mr. Genie-what is a
                 meta-wish? Genie: It is simply a wish about wishes. I am not allowed
                 to grant meta-wishes. It is only within my purview to grant plain
                 ordinary wishes, such as wishing for ten bottles of beer, to have Helen
                 of Troy on a blanket, or to have an all-expenses-paid weekend for two
                 at the Copacabana. You know-simple things like that. But meta-
                 wishes I cannot grant. GOD won't permit me to.
          Achilles: GOD? Who is GOD? And why won't he let you grant meta-wishes?
                 That seems like such a puny thing compared to the others you
                 mentioned.




Little Harmonic Labyrinth                                                           118

          Genie: Well, it’s a complicated matter, you see. Why don’t you just go ahead
                 and make your three wishes? Or at least make one of them. I don't
                 have all I time in the world, you know ...
          Achilles: Oh, I feel so rotten. I was REALLY HOPING wish for a hundred
                 wishes ...
          Genie: Gee, I hate to see anybody so disappointed that. And besides, meta-
                 wishes are my favorite k of wish. Let me just see if there isn't anything
                 I do about this. This'll just take one moment

             (The Genie removes from the wispy folds of his robe an object which looks
             just like the copper Lamp he had put away, except that this one is made of
             silver; and where the previous one had 'L' etched on it, this one has 'ML' in
             smaller letters, so as to cover the same area.)
          I
          Achilles: And what is that?
          Genie: This is my Meta-Lamp ...

             (He rubs the Meta-Lamp, and a huge puff of smoke appears. In the billows
             of smoke, they can all make out a ghostly form towering above them.)

          Meta-Genie: I am the Meta-Genie. You summoned me, 0 Genie? What is
                 your wish?
          Genie: I have a special wish to make of you, 0 Djinn and of GOD. I wish for
                 permission for tempos suspension of all type-restrictions on wishes,
                 for duration of one Typeless Wish. Could you ph grant this wish for
                 me?
          Meta-Genie: I'll have to send it through Channels, of course. One half a
                 moment, please

             (And, twice as quickly as the Genie did, this Meta-Genie removes from the
             wispy folds of her robe an object which looks just like the silver Meta-
             Lamp, except that it is made of gold; and where the previous one had 'ML'
             etched on it, this one has 'MML' in smaller letters, so as to cover the same
             area.)

          Achilles (his voice an octave higher than before): And what is that? Meta-
                 Genie: This is my Meta-Meta-Lamp. . .

             (She rubs the Meta-Meta-Lamp, and a hugs puff of smoke appears. In the
             billows o smoke, they can all make out a ghostly fore towering above
             them.)




Little Harmonic Labyrinth                                                             119

            Meta-Meta-Genie: I am the MetaMeta-Genie. You summoned me,
          0 Meta-Genie? What is your wish?
          Meta-Genie: I have a special wish to make of you, 0 Djinn, and of GOD. I
                 wish for permission for temporary suspension of all type-restrictions
                 on wishes, for the duration of one Typeless Wish. Could you please
                 grant this wish for me?
                   Meta-Meta-Genie: I'll have to send it through Channels, of course.
                            One quarter of a moment, please.

                       (And, twice as quickly as the Meta-Genie did, this MetaMeta-
                       Genie removes from the folds of his robe an object which looks
                       just like the gold MetaLamp, except that it is made of ...)
                       .
                           .
                               .
                                   .
                                           .
                                               . .{GOD}
                                           .
                                       .
                               .
                       .
                       ( ... swirls back into the MetaMeta-Meta-Lamp, which the Meta-
                       Meta-Genie then folds back into his robe, half as quickly as the
                       Meta-Meta-Meta-Genie did.)

                       Your wish is granted, 0 MetaGenie.

                    Meta-Genie: Thank you, 0 Djinn, and GOD.

                  (And the Meta-Meta-Genie, as all the higher ones before him, swirls
                  back into the Meta-Meta-Lamp, which the Meta-Genie then folds back
                  into her robe, half as quickly as the Meta-Meta-Genie did.)

                 Your wish is granted, 0 Genie.
          Genie: Thank you, 0 Djinn, and GOD.

              (And the Meta-Genie, as all the higher ones before her,




Little Harmonic Labyrinth                                                          120

          swirls back into the Meta-Lamp, which the Genie folds back into his robe,
          half as quickly as the M Genie did.)

             Your wish is granted, Achilles.

             (And one precise moment has elapsed since he "This will just take one
             moment.")

          Achilles: Thank you, 0 Djinn, and GOD.
          Genie: I am pleased to report, Achilles, that you r have exactly one (1)
                 Typeless Wish-that is to sa wish, or a meta-wish, or a meta-meta-wish,
                 as many "meta"'s as you wish-even infinitely many (if wish).
          Achilles: Oh, thank you so very much, Genie. But curiosity is provoked.
                 Before I make my wish, would you mind telling me who-or what-
                 GOD is?
          Genie: Not at all. "GOD" is an acronym which stands "GOD Over Djinn".
                 The word "Djinn" is used designate Genies, Meta-Genies, Meta-Meta-
                 Gen etc. It is a Typeless word.
          Achilles: But-but-how can "GOD" be a word in own acronym? That doesn't
                 make any sense!
          Genie: Oh, aren't you acquainted with recursive acronyms? I thought
                 everybody knew about them. \ see, "GOD" stands for "GOD Over
                 Djinn"-which can be expanded as "GOD Over Djinn, O, Djinn"-and
                 that can, in turn, be expanded to "G( Over Djinn, Over Djinn, Over
                 Djinn"-which can its turn, be further expanded ... You can go as as
                 you like.
          Achilles: But I'll never finish!
          Genie: Of course not. You can never totally expand GOD.
          Achilles: Hmm ... That's puzzling. What did you me when you said to the
                 Meta-Genie, "I have a sped wish to make of you, 0 Djinn, and of
                 GOD"?
          Genie: I wanted not only to make a request of Meta-Genie, but also of all the
                 Djinns over her. 'I recursive acronym method accomplishes this qL
                 naturally. You see, when the Meta-Genie received my request, she
                 then had to pass it upwards to I GOD. So she forwarded a similar
                 message to I Meta-Meta-Genie, who then did likewise to t Meta-Meta-
                 Meta-Genie ... Ascending the chain this way transmits the message to
                 GOD.




Little Harmonic Labyrinth                                                          121

          Achilles: I see. You mean GOD sits up at the top of the ladder of djinns?
          Genie: No, no, no! There is nothing "at the top", for there is no top. That is
                 why GOD is a recursive acronym. GOD is not some ultimate djinn;
                 GOD is the tower of djinns above any given djinn.
          Tortoise: It seems to me that each and every djinn would have a different
                 concept of what GOD is, then, since to any djinn, GOD is the set of
                 djinns above him or her, and no two djinns share that set.
          Genie: You're absolutely right-and since I am the lowest djinn of all, my
                 notion of GOD is the most exalted one. I pity the higher djinns, who
                 fancy themselves somehow closer to GOD. What blasphemy!
          Achilles: By gum, it must have taken genies to invent GOD.
          Tortoise: Do you really believe all this stuff about GOD, Achilles?
          Achilles: Why certainly, I do. Are you atheistic, Mr. T? Or are you agnostic?
          Tortoise: I don't think I'm agnostic. Maybe I'm metaagnostic.
          Achilles: Whaaat? I don't follow you at all.
          Tortoise: Let's see . . . If I were meta-agnostic, I'd be confused over whether
                 I'm agnostic or not-but I'm not quite sure if I feel THAT way; hence I
                 must be meta-meta-agnostic (I guess). Oh, well. Tell me, Genie, does
                 any djinn ever make a mistake, and garble up a message moving up or
                 down the chain?
          Genie: This does happen; it is the most common cause for Typeless Wishes
                 not being granted. You see, the chances are infinitesimal, that a
                 garbling will occur at any PARTICULAR link in the chain-but when
                 you put an infinite number of them in a row, it becomes virtually
                 certain that a garbling will occur SOMEWHERE. In fact, strange as it
                 seems, an infinite number of garblings usually occur, although they
                 are very sparsely distributed in the chain.
          Achilles: Then it seems a miracle that any Typeless Wish ever gets carried
                 out.
          Genie: Not really. Most garblings are inconsequential, and many garblings
                 tend to cancel each other out. But occasionally-in fact, rather seldom-
                 the nonfulfillment of a Typeless Wish can be traced back to a single
                 unfortunate djinn's garbling. When this happens, the guilty djinn is
                 forced to run an infinite




Little Harmonic Labyrinth                                                            122

          Gauntlet and get paddled on his or her rump, by GOD. It's good fun for the
                 paddlers, and q harmless for the paddlee. You might be amused by the
                 sight.
          Achilles: I would love to see that! But it only happens when a Typeless Wish
                 goes ungranted?
          Genie: That's right.
          Achilles: Hmm ... That gives me an idea for my w Tortoise: Oh, really? What
                 is it? Achilles: I wish my wish would not be granted!

             (At that moment, an event-or is "event" the word for it? --takes place which
             cannot be described, and hence no attempt will be made to describe it.)

          Achilles: What on earth does that cryptic comment mean?
          Tortoise: It refers to the Typeless Wish Achilles made.
          Achilles: But he hadn't yet made it.
          Tortoise: Yes, he had. He said, "I wish my wish would not be
          granted", and the Genie took THAT to be his wish.

             (At that moment, some footsteps are heard coming down the hallway in
             their direction.)

          Achilles: Oh, my! That sounds ominous.

              (The footsteps stop; then they turn around and fade away.)

          Tortoise: Whew!
          Achilles: But does the story go on, let's see. or is that the end? Turn the page
                 and let’s see.

             (The Tortoise turns the page of "Djinn and Tonic", where they find that the
             story goes on ...)

               Achilles: Hey! What happened? Where is my Genie: lamp? My cup of
                      espresso? What happened to young friends from the Convex and
                      Concave worlds? What are all those little lizards doing hi
               Tortoise: I'm afraid our context got restored incorrectly Achilles.
               Achilles: What on earth does that cryptic comment mean?
               Tortoise: I refer to the Typeless Wish you made.
               Achilles: But I hadn't yet made it.
               Tortoise: Yes, you had. You said, "I wish my wish would not be
                      granted", and the Genie took THAT to be your wish.
               Achilles: Oh, my! That sounds ominous.
               Tortoise; It spells PARADOX. For that Typeless wish to be




Little Harmonic Labyrinth                                                             123

                    granted, it had to be denied – yet not to grant it would be to grant
                    it.
               Achilles: So what happened? Did the earth come to a standstill? Did the
                      universe cave in?
               Tortoise: No. The System crashed. Achilles: What does that mean?
               Tortoise: It means that you and I, Achilles, were suddenly and
                      instantaneously transported to Tumbolia. Achilles: To where?
               Tortoise: Tumbolia: the land of dead hiccups and extinguished light
                      bulbs. It's a sort of waiting room, where dormant software waits
                      for its host hardware to come back up. No telling how long the
                      System was down, and we were in Tumbolia. It could have been
                      moments, hours, days-even years.
               Achilles: I don't know what software is, and I don't know what hardware
                      is. But I do know that I didn't get to make my wishes! I want my
                      Genie back!
               Tortoise: I'm sorry, Achilles-you blew it. You crashed the System, and
                      you should thank your lucky stars that we're back at all. Things
                      could have come out a lot worse. But I have no idea where we
                      are.
               Achilles: I recognize it now-we're inside another of Escher's pictures.
                      This time it's Reptiles.
               Tortoise: Aha! The System tried to save as much of our context as it
                      could before it crashed, and it got as far as recording that it was
                      an Escher picture with lizards before it went down. That's
                      commendable.
               Achilles: And look-isn't that our phial of poppingtonic over there on the
                      table, next to the cycle of lizards?
               Tortoise: It certainly is, Achilles. I must say, we are very lucky indeed.
                      The System was very kind to us, in giving us back our popping-
                      tonic-it's precious stuff!
               Achilles: I'll say! Now we can pop back out of the Escher world, into my
                      house.
               Tortoise: There are a couple of books on the desk, next to the tonic. I
                      wonder what they are. (He picks up the smaller one, which is
                      open to a random page.) This looks like a moderately
                      provocative book.
               Achilles: Oh, really? What is its title?
               Tortoise: Provocative Adventures of the Tortoise and Achilles Taking
                      Place in Sundry Parts of the Globe. It sounds like an interesting
                      book to read out of.




Little Harmonic Labyrinth                                                            124

                      FIGURE 24. Reptiles, by M. C. Escher (lithograph, 1943).

               Achilles: Well, You can read it if you want, but as for I'm not going to
                      take any chances with t popping-tonic-one of the lizards might
                      knock it off the table, so I'm going to get it right now!

                  (He dashes over to the table and reaches for the popping-tonic, but in
                  his haste he somehow bumps the flask of tonic, and it tumbles off the
                  desk and begins rolling.)

               Oh, no! Mr. T-look! I accidentally knocked tonic onto the floor, and it's
                      rolling toward towards-the stairwell! Quick-before it falls!

               (The Tortoise, however, is completely wrapped up in the thin volume
                     which he has in his hands.) Achilles: Well, You can read it if you
                     want, but as for I'm not going to take any chances with t popping-
                     tonic-one of the lizards might knock it off the table, so I'm going
                     to get it right




Little Harmonic Labyrinth                                                           125

               Tortoise (muttering): Eh? This story looks fascinating.
               Achilles: Mr. T, Mr. T, help! Help catch the tonic-flask!
               Tortoise: What's all the fuss about?
               Achilles: The tonic-flask-I knocked it down from the desk, and now it's
                      rolling and

                  (At that instant it reaches the brink of the stairwell, and plummets
                  over ... )

                 Oh no! What can we do? Mr. Tortoise-aren't you alarmed? We're
                 losing our tonic! It's just fallen down the stairwell! There's only one
                 thing to do! We'll have to go down one story!
               Tortoise: Go down one story? My pleasure. Won't you join me?

                  (He begins to read aloud, and Achilles, pulled in two directions at
                  once, finally stays, taking the role of the Tortoise.)

               Achilles: It's very dark here, Mr. T. I can't see a thing. Oof! I bumped
                      into a wall. Watch out!
               Tortoise: Here-I have a couple of walking sticks. Why don't you take one
                      of them? You can hold it out in front of you so that you don't
                      bang into things.
               Achilles: Good idea. (He takes the stick.) Do you get the sense that this
                      path is curving gently to the left as we walk? Tortoise: Very
                      slightly, yes.
               Achilles: I wonder where we are. And whether we'll ever see the light of
                      day again. I wish I'd never listened to you, when you suggested I
                      swallow some of that "DRINK ME" stuff.
               Tortoise: I assure you, it's quite harmless. I've done it scads of times, and
                      not a once have I ever regretted it. Relax and enjoy being small.
               Achilles: Being small? What is it you've done to me, Mr. T?
               Tortoise: Now don't go blaming me. You did it of your own free will.
                      Achilles: Have you made me shrink? So that this labyrinth we're
                      in is actually some teeny thing that someone could STEP on?




Little Harmonic Labyrinth                                                               126

               FIGURE 25. Cretan Labyrinth (Italian engraving; School of
                    Finiguerra). [From N Matthews, Mazes and Labyrinths: Their
                    History and Development (New York: Dover Publications, 1970).

               Tortoise: Labyrinth? Labyrinth? Could it Are we in the notorious Little
                      Harmonic Labyrinth of the dreaded Majotaur?
               Achilles: Yiikes! What is that?
               Tortoise: They say-although I person never believed it myself-that an I
                      Majotaur has created a tiny labyrinth sits in a pit in the middle of
                      it, waiting innocent victims to get lost in its fears complexity.
                      Then, when they wander and dazed into the center, he laughs and
                      laughs at them-so hard, that he laughs them to death!
               Achilles: Oh, no!
               Tortoise: But it's only a myth. Courage, Achilles.

                      (And the dauntless pair trudge on.)

               Achilles: Feel these walls. They're like o gated tin sheets, or something.
                      But the corrugations have different sizes.




Little Harmonic Labyrinth                                                             127

                  (To emphasize his point, he sticks out his walking stick against the
                  wall surface as he walks. As the stick bounces back and forth against
                  the corrugations, strange noises echo up and down the long curved
                  corridor they are in.)

               Tortoise (alarmed): What was THAT?
               Achilles: Oh, just me, rubbing my walking stick against the wall.
               Tortoise: Whew! I thought for a moment it was the bellowing of the
                      ferocious Majotaur! Achilles: I thought you said it was all a
                      myth.
               Tortoise: Of course it is. Nothing to be afraid of.

                  (Achilles puts his walking stick back against the wall, and continues
                  walking. As he does so, some musical sounds are heard, coming from
                  the point where his stick is scraping the wall.)

               Tortoise: Uh-oh. I have a bad feeling, Achilles.
               That Labyrinth may not be a myth, after all. Achilles: Wait a minute.
                      What makes you change your mind all of a sudden? Tortoise: Do
                      you hear that music?

                  (To hear more clearly, Achilles lowers the stick, and the strains of
                  melody cease.)

                  Hey! Put that back! I want to hear the end of this piece!

                  (Confused, Achilles obeys, and the music resumes.)

                 Thank you. Now as I was about to say, I have just figured out where
               we are.
               Achilles: Really? Where are we?

               Tortoise: We are walking down a spiral groove of a record in its jacket.
                      Your stick scraping against the strange shapes in the wall acts
                      like a needle running down the groove, allowing us to hear the
                      music.
               Achilles: Oh, no, oh, no ...
               Tortoise: What? Aren't you overjoyed? Have you ever had the chance to
                      be in such intimate contact with music before?




Little Harmonic Labyrinth                                                          128

               Achzltes: How am I ever going to win footraces against full-sized people
                      when I am smaller than a flea, Mr. Tortoise?
               Tortoise: Oh, is that all that's bothering you That's nothing to fret abopt,
                      Achilles.
               Achilles: The way you talk, I get the impression that you never worry at
                      all.
               Tortoise: I don't know. But one thing for certain is that I don't worry
                      about being small. Especially not when faced with the awful
                      danger of the dreaded Majotaur!
               Achilles: Horrors! Are you telling me
               Tortoise: I'm afraid so, Achilles. The music gave it away.
               Achilles: How could it do that?
               Tortoise: Very simple. When I heard melody B-A-C-H in the top voice,
                      I immediately realized that the grooves we're walking through
                      could only be Little Harmonic Labyrinth, one of Bach's er known
                      organ pieces. It is so named cause of its dizzyingly frequent
                      modulations.
               Achilles: Wh-what are they?
               Tortoise: Well, you know that most music pieces are written in a key, or
                      tonality, as C major, which is the key of this o;
               Achilles: I had heard the term before. Do that mean that C is the note
                      you want to on?
               Tortoise: Yes, C acts like a home base, in a Actually, the usual word is
                      "tonic".
               Achilles: Does one then stray away from tonic with the aim of eventually
                      returning
               Tortoise: That's right. As the piece develops ambiguous chords and
                      melodies are t which lead away from the tonic. Little by little,
                      tension builds up-you feel at creasing desire to return home, to
                      hear the tonic.
               Achilles: Is that why, at the end of a pie always feel so satisfied, as if I
                      had waiting my whole life to hear the ton
               Tortoise: Exactly. The composer has uses knowledge of harmonic
                      progressions to




Little Harmonic Labyrinth                                                              129

                    manipulate your emotions, and to build up hopes in you to hear
                    that tonic.
               Achilles: But you were going to tell me about modulations.
               Tortoise: Oh, yes. One very important thing a composer can do is to
                      "modulate" partway through a piece, which means that he sets up
                      a temporary goal other than resolution into the tonic.
               Achilles: I see ... I think. Do you mean that some sequence of chords
                      shifts the harmonic tension somehow so that I actually desire to
                      resolve in a new key?
               Tortoise: Right. This makes the situation more complex, for although in
                      the short term you want to resolve in the new key, all the while at
                      the back of your mind you retain the longing to hit that original
                      goal-in this case, C major. And when the subsidiary goal is
                      reached, there is
               Achilles (suddenly gesturing enthusiastically): Oh, listen to the gorgeous
                      upward-swooping chords which mark the end of this Little
                      Harmonic Labyrinth!
               Tortoise: No, Achilles, this isn't the end. It's merely
               Achilles: Sure it is! Wow! What a powerful, strong ending! What a sense
                      of relief! That's some resolution! Gee!

                  (And sure enough, at that moment the music stops, as they emerge into
                  an open area with no walls.)

               You see, it Is over. What did I tell you? Tortoise: Something is very
                       wrong. This record
               is a disgrace to the world of music. Achilles: What do you mean?
               Tortoise: It was exactly what I was telling you about. Here Bach had
                       modulated from C into G, setting up a secondary goal of hearing
                       G. This means that you experience two tensions at once-waiting
                       for resolution into G, but also keeping in mind that ultimate
                       desire-to resolve triumphantly into C Major.
               Achilles: Why should you have to keep any




Little Harmonic Labyrinth                                                            130

                    thing in mind when listening to a piece of music? Is music only an
                    intellectual exercise?
               Tortoise: No, of course not. Some music is highly intellectual, but most
                      music is not. And most of the time your ear or br the
                      "calculation" for you, and lets your emotions know what they
                      want to hear, don't have to think about it consciously in this
                      piece, Bach was playing tricks hoping to lead you astray. And in
                      your case Achilles, he succeeded.
               Achilles: Are you telling me that I responded to a resolution in a
                      subsidiary key?
               Tortoise: That's right.
               Achilles: It still sounded like an ending to me
               Tortoise: Bach intentionally made it sot way. You just fell into his trap.
                      It was deliberately contrived to sound like an ending but if you
                      follow the harmonic progression carefully, you will see that it is
                      in the wrong key. Apparently not just you but this miserable
                      record company fell for the same trick-and they truncated the
                      piece early.
               Achilles: What a dirty trick Bach played
               Tortoise: That is his whole game-to m lose your way in his Labyrinth! 'l
                      Majotaur is in cahoots with Bach, And if you don't watch out, he
                      i laugh you to death-and perhaps n with you!
               Achilles: Oh, let us hurry up and get here! Quick! Let's run backwards
                      grooves, and escape on the outside record before the Evil
                      Majotaur finds us.
               Tortoise: Heavens, no! My sensibility is delicate to handle the bizarre the
                      gressions which occur when time versed.
               Achilles: Oh, Mr. T, how will we ever get out of here, if we can't just
                      retrace our steps
               Tortoise: That's a very good question.

                  (A little desperately, Achilles starts runt about aimlessly in the dark.
                  Suddenly t is a slight gasp, and then a "thud".)




Little Harmonic Labyrinth                                                             131

               Achilles-are you all right?
               Achilles: Just a bit shaken up but otherwise fine. I fell into some big
                      hole.
               Tortoise: You've fallen into the pit of the Evil Majotaur! Here, I'll come
                      help you out. We've got to move fast!
               Achilles: Careful, Mr. T-I don't want You to fall in here, too ...
               Tortoise: Don't fret, Achilles. Everything will be all --

                     (Suddenly, there is a slight gasp, and then a "thud".)

               Achilles: Mr. T-you fell in, too! Are you all right?
               Tortoise: Only my pride is hurt-otherwise I'm fine.
               Achilles: Now we're in a pretty pickle, aren't we?

                      (Suddenly, a giant, booming laugh is heard, alarmingly close to
                     them.)

               Tortoise: Watch out, Achilles! This is no laughing matter.
               Majotaur: Hee hee hee! Ho ho! Haw haw haw!
               Achilles: I'm starting to feel weak, Mr. T ...
               Tortoise: Try to pay no attention to his laugh,
               Achilles. That's your only hope.
               Achilles: I'll do my best. If only my stomach weren't empty!
               Tortoise: Say, am I smelling things, or is there a bowl of hot buttered
                      popcorn around here? Achilles: I smell it, too. Where is it coming
                      from?
               Tortoise: Over here, I think. Oh! I just ran into a big bowl of the stuff.
                      Yes, indeed-it seems to be a bowl of popcorn!
               Achilles: Oh, boy-popcorn! I'm going to munch my head off!

        Tortoise: Let's just hope it isn't pushcorn! Pushcorn and popcorn are
               extraordinarily difficult to tell apart.

               Achilles: What's this about Pushkin?
               Tortoise: I didn't say a thing. You must be hearing things.
               Achilles: Go-golly! I hope not. Well, let's dig in!




Little Harmonic Labyrinth                                                            132

                  (And the two Jriends begin muncnai popcorn (or pushcorn?)-and t
                  once POP! I guess it was popcorn; all.)

               Tortoise: What an amusing story. Did you en
               Achilles: Mildly. Only I wonder whether the' out of that Evil Majotaur's
                      pit or r Achilles-he wanted to be full-sized again
               Tortoise: Don't worry-they're out, and he is again. That's what the "POP"
                      was all abo
               Achilles: Oh, I couldn't tell. Well, now I REAL: find that bottle of tonic.
                      For some reason, burning. And nothing would taste bett drink of
                      popping-tonic.
               Tortoise: That stuff is renowned for its thirst powers. Why, in some
                      places people very crazy over it. At the turn of the century the
                      Schonberg food factory stopped ma] and started making cereal
                      instead. You cai the uproar that caused.
               Achilles: I have an inkling. But let's go look fo Hey just a moment.
                      Those lizards on the you see anything funny about them?
               Tortoise: Umm ... not particularly. What do you see of such great
                      interest?
               Achilles: Don't you see it? They're emerging flat picture without
                      drinking any pop] How are they able to do that?
               Tortoise: Oh, didn't I tell you? You can ge picture by moving
                      perpendicularly to it you have no popping-tonic. The little li
                      learned to climb UP when they want to ge two-dimensional
                      sketchbook world.
               Achilles: Could we do the same thing to get Escher picture we're in?
               Tortoise: Of course! We just need to go UP one story. you want to try it?
               Achilles: Anything to get back to my house! I all these provocative
                      adventures.
               Tortoise: Follow me, then, up this way.

                     (And they go up one story.)

               Achilles: It's good to be back. But something seems wrong. This isn't my
                      house! This is YOUR house, Mr. Tortoise
               Tortoise: Well, so it is-and am I glad for that! I wasn’t looking




Little Harmonic Labyrinth                                                             133

                 forward one whit to the long walk back from your house. I am bushed,
                 and doubt if I could have made it.
               Achilles: I don't mind walking home, so I guess it's lucky we ended up
                      here, after all.
               Tortoise: I'll say! This certainly is a piece of Good Fortune!




Little Harmonic Labyrinth                                                        134

                          Recursive Structures
                             and Processes

                                What Is Recursion?

WHAT IS RECURSION? It is what was illustrated in the Dialogue Little Harmonic
Labyrinth: nesting, and variations on nesting. The concept is very general. (Stories inside
stories, movies inside movies, paintings inside paintings, Russian dolls inside Russian
dolls (even parenthetical comments in. side parenthetical comments!)-these are just a few
of the charms of recursion.) However, you should he aware that the meaning of
"recursive' in this Chapter is only faintly related to its meaning in Chapter 111. The
relation should be clear by the end of this Chapter.
        Sometimes recursion seems to brush paradox very closely. For example, there are
recursive definitions. Such a definition may give the casual viewer the impression that
something is being defined in terms of itself. That would be circular and lead to infinite
regress, if not to paradox proper. Actually, a recursive definition (when properly
formulated) never leads to infinite regress or paradox. This is because a recursive
definition never defines something in terms of itself, but always in terms of simpler
versions of itself. What I mean by this will become clearer shortly, when ' show some
examples of recursive definitions.
        One of the most common ways in which recursion appears in daily life is when
you postpone completing a task in favor of a simpler task, often o the same type. Here is
a good example. An executive has a fancy telephone and receives many calls on it. He is
talking to A when B calls. To A he say,, "Would you mind holding for a moment?" Of
course he doesn't really car if A minds; he just pushes a button, and switches to B. Now C
calls. The same deferment happens to B. This could go on indefinitely, but let us not get
too bogged down in our enthusiasm. So let's say the call with C terminates. Then our
executive "pops" back up to B, and continues. Meanwhile A is sitting at the other end of
the line, drumming his fingernails again some table, and listening to some horrible
Muzak piped through the phone lines to placate him ... Now the easiest case is if the call
with B simply terminates, and the executive returns to A finally. But it could happen that
after the conversation with B is resumed, a new caller-D-calls. B is once again pushed
onto the stack of waiting callers, and D is taken care of. Aft D is done, back to B, then
back to A. This executive is hopelessly mechanical, to be sure-but we are illustrating
recursion in its most precise form




Recursive Structures and Processes                                                     135

                           Pushing, Popping, and Stacks
In the preceding example, I have introduced some basic terminology of recursion-at least
as seen through the eyes of computer scientists. The terms are push, pop, and stack (or
push-down stack, to be precise) and they are all related. They were introduced in the late
1950's as part of IPL, one of the first languages for Artificial Intelligence. You have
already encountered "push" and "pop" in the Dialogue. But I will spell things out
anyway. To push means to suspend operations on the task you're currently working on,
without forgetting where you are-and to take up a new task. The new task is usually said
to be "on a lower level" than the earlier task. To pop is the reverse-it means to close
operations on one level, and to resume operations exactly where you left off, one level
higher.
         But how do you remember exactly where you were on each different level? The
answer is, you store the relevant information in a stack. So a stack is just a table telling
you such things as (1) where you were in each unfinished task (jargon: the "return
address"), (2) what the relevant facts to know were at the points of interruption (jargon:
the "variable bindings"). When you pop back up to resume some task, it is the stack
which restores your context, so you don't feel lost. In the telephone-call example, the
stack tells you who is waiting on each different level, and where you were in the
conversation when it was interrupted.
         By the way, the terms "push", "pop", and "stack" all come from the visual image
of cafeteria trays in a stack. There is usually some sort of spring underneath which tends
to keep the topmost tray at a constant height, more or less. So when you push a tray onto
the stack, it sinks a little-and when you remove a tray from the stack, the stack pops up a
little.
         One more example from daily life. When you listen to a news report on the radio,
oftentimes it happens that they switch you to some foreign correspondent. "We now
switch you to Sally Swumpley in Peafog, England." Now Sally has got a tape of some
local reporter interviewing someone, so after giving a bit of background, she plays it. "I'm
Nigel Cadwallader, here on scene just outside of Peafog, where the great robbery took
place, and I'm talking with ..." Now you are three levels down. It may turn out that the
interviewee also plays a tape of some conversation. It is not too uncommon to go down
three levels in real news reports, and surprisingly enough, we scarcely have any
awareness of the suspension. It is all kept track of quite easily by our subconscious mind.
Probably the reason it is so easy is that each level is extremely different in flavor from
each other level. If they were all similar, we would get confused in no time flat.
         An example of a more complex recursion is, of course, our Dialogue. There,
Achilles and the Tortoise appeared on all the different levels. Sometimes they were
reading a story in which they appeared as characters. That is when your mind may get a
little hazy on what's going on, and you have to concentrate carefully to get things straight.
"Let's see, the real Achilles and Tortoise are still up there in Goodfortune's helicopter, but
the




Recursive Structures and Processes                                                        136

secondary ones are in some Escher picture-and then they found this book and are reading
in it, so it's the tertiary Achilles and Tortoise who wandering around inside the grooves
of the Little Harmonic Labyrinth. wait a minute-I left out one level somewhere ..." You
have to ha conscious mental stack like this in order to keep track of the recursion the
Dialogue. (See Fig. 26.)




FIGURE 26. Diagram of the structure of the Dialogue Little Harmonic Labyrinth
Vertical descents are "pushes"; rises ore "pops". Notice the similarity of this diagram to
indentation pattern of the Dialogue. From the diagram it is clear that the initial tension
Goodfortune's threat-never was resolved; Achilles and the Tortoise were just left
dangling the sky. Some readers might agonize over this unpopped push, while others
might not ba eyelash. In the story, Bach's musical labyrinth likewise was cut off too soon-
but Achilles d even notice anything funny. Only the Tortoise was aware of the more
global dangling tension

                                   Stacks in Music
        While we're talking about the Little Harmonic Labyrinth, we should discuss
something which is hinted at, if not stated explicitly in the Dialogue: that hear music
recursively-in particular, that we maintain a mental stack of keys, and that each new
modulation pushes a new key onto the stack. implication is further that we want to hear
that sequence of keys retrace reverse order-popping the pushed keys off the stack, one by
one, until the tonic is reached. This is an exaggeration. There is a grain of truth to it
however.
        Any reasonably musical person automatically maintains a shallow with two keys.
In that "short stack", the true tonic key is held and also most immediate "pseudotonic"
(the key the composer is pretending t in). In other words, the most global key and the
most local key. That the listener knows when the true tonic is regained, and feels a strong
s of "relief". The listener can also distinguish (unlike Achilles) between a local easing of
tension-for example a resolution into the pseudotonic --




Recursive Structures and Processes                                                      137

and a global resolution. In fact, a pseudoresolution should heighten the global tension,
not relieve it, because it is a piece of irony-just like Achilles' rescue from his perilous
perch on the swinging lamp, when all the while you know he and the Tortoise are really
awaiting their dire fates at the knife of Monsieur Goodfortune.
        Since tension and resolution are the heart and soul of music, there are many, many
examples. But let us just look at a couple in Bach. Bach wrote many pieces in an
"AABB" form-that is, where there are two halves, and each one is repeated. Let's take the
gigue from the French Suite no. 5, which is quite typical of the form. Its tonic key is G,
and we hear a gay dancing melody which establishes the key of G strongly. Soon,
however, a modulation in the A-section leads to the closely related key of D (the
dominant). When the A-section ends, we are in the key of D. In fact, it sounds as if the
piece has ended in the key of D! (Or at least it might sound that way to Achilles.) But
then a strange thing happens-we abruptly jump back to the beginning, back to G, and
rehear the same transition into D. But then a strange thing happens-we abruptly jump
back to the beginning, back to G, and rehear the same transition into D.
        Then comes the B-section. With the inversion of the theme for our melody, we
begin in D as if that had always been the tonic-but we modulate back to G after all, which
means that we pop back into the tonic, and the B-section ends properly. Then that funny
repetition takes place, jerking us without warning back into D, and letting us return to G
once more. Then that funny repetition takes place, jerking us without warning
        back into D, and letting us return to G once more.
        The psychological effect of all this key shifting-some jerky, some smooth-is very
difficult to describe. It is part of the magic of music that we can automatically make sense
of these shifts. Or perhaps it is the magic of Bach that he can write pieces with this kind
of structure which have such a natural grace to them that we are not aware of exactly
what is happening.
        The original Little Harmonic Labyrinth is a piece by Bach in which he tries to
lose you in a labyrinth of quick key changes. Pretty soon you are so disoriented that you
don't have any sense of direction left-you don't know where the true tonic is, unless you
have perfect pitch, or like Theseus, have a friend like Ariadne who gives you a thread that
allows you to retrace your steps. In this case, the thread would be a written score. This
piece-another example is the Endlessly Rising Canon-goes to show that, as music
listeners, we don't have very reliable deep stacks.


                               Recursion in Language
Our mental stacking power is perhaps slightly stronger in language. The grammatical
structure of all languages involves setting up quite elaborate push-down stacks, though, to
be sure, the difficulty of understanding a sentence increases sharply with the number of
pushes onto the stack. The proverbial German phenomenon of the "verb-at-the-end",
about which




Recursive Structures and Processes                                                      138

       Droll tales of absentminded professors who would begin a sentence, ramble on for
an entire lecture, and then finish up by rattling off a string of verbs by which their
audience, for whom the stack had long since lost its coherence, would be totally
nonplussed, are told, is an excellent example of linguistic pushing and popping. The
confusion among the audience out-of-order popping from the stack onto which the
professor's verbs been pushed, is amusing to imagine, could engender. But in normal ken
German, such deep stacks almost never occur-in fact, native speaker of German often
unconsciously violate certain conventions which force verb to go to the end, in order to
avoid the mental effort of keeping track of the stack. Every language has constructions
which involve stacks, though usually of a less spectacular nature than German. But there
are always of rephrasing sentences so that the depth of stacking is minimal.

                          Recursive Transition Networks
The syntactical structure of sentences affords a good place to present a of describing
recursive structures and processes: the Recursive Transition Network (RTN). An RTN is
a diagram showing various paths which can be followed to accomplish a particular task.
Each path consists of a number of nodes, or little boxes with words in them, joined by
arcs, or lines with arrows. The overall name for the RTN is written separately at the left,
and the and last nodes have the words begin and end in them. All the other nodes contain
either very short explicit directions to perform, or else name other RTN's. Each time you
hit a node, you are to carry out the direct inside it, or to jump to the RTN named inside it,
and carry it out.
        Let's take a sample RTN, called ORNATE NOUN, which tells how to construct
a certain type of English noun phrase. (See Fig. 27a.) If traverse ORNATE NOUN
purely horizontally, we begin', then we create ARTICLE, an ADJECTIVE, and a
NOUN, then we end. For instance, "the shampoo" or "a thankless brunch". But the arcs
show other possibilities such as skipping the article, or repeating the adjective. Thus we
co construct "milk", or "big red blue green sneezes", etc.
        When you hit the node NOUN, you are asking the unknown black I called NOUN
to fetch any noun for you from its storehouse of nouns. This is known as a procedure
call, in computer science terminology. It means you temporarily give control to a
procedure (here, NOUN) which (1) does thing (produces a noun) and then (2) hands
control back to you. In above RTN, there are calls on three such procedures: ARTICLE,
ADJECTIVE and NOUN. Now the RTN ORNATE NOUN could itself be called from
so other RTN-for instance an RTN called SENTENCE. In this case, ORNATE NOUN
would produce a phrase such as "the silly shampoo" and d return to the place inside
SENTENCE from which it had been called. I quite reminiscent of the way in which you
resume where you left off nested telephone calls or nested news reports.
        However, despite calling this a "recursive transition network", we have




Recursive Structures and Processes                                                       139

FIGURE 27. Recursive Transition Networks for ORNATE NOUN and FANCY NOUN.

not exhibited any true recursion so far. Things get recursive-and seemingly circular-when
you go to an RTN such as the one in Figure 27b, for FANCY NOUN. As you can see,
every possible pathway in FANCY NOUN involves a call on ORNATE NOUN, so there
is no way to avoid getting a noun of some sort or other. And it is possible to be no more
ornate than that, coming out merely with "milk" or "big red blue green sneezes". But
three of the pathways involve recursive calls on FANCY NOUN itself. It certainly looks
as if something is being defined in terms of itself. Is that what is happening, or not?
        The answer is "yes, but benignly". Suppose that, in the procedure SENTENCE,
there is a node which calls FANCY NOUN, and we hit that node. This means that we
commit to memory (viz., the stack) the location of that node inside SENTENCE, so we'll
know where to return to-then we transfer our attention to the procedure FANCY NOUN.
Now we must choose a pathway to take, in order to generate a FANCY NOUN. Suppose
we choose the lower of the upper pathways-the one whose calling sequence goes:

       ORNATE NOUN; RELATIVE PRONOUN; FANCY NOUN; VERB.




Recursive Structures and Processes                                                   140

So we spit out an ORNATE NOUN: "the strange bagels"; a RELATIVE NOUN:
"that"; and now we are suddenly asked for a FANCY NOUN. B are in the middle of
FANCY NOUN! Yes, but remember our executive was in the middle of one phone call
when he got another one. He n stored the old phone call's status on a stack, and began the
new one nothing were unusual. So we shall do the same.
        We first write down in our stack the node we are at in the outer call on FANCY
NOUN, so that we have a "return address"; then we jump t beginning of FANCY NOUN
as if nothing were unusual. Now we h~ choose a pathway again. For variety's sake, let's
choose the lower pat] ORNATE NOUN; PREPOSITION; FANCY NOUN. That
means we produce an ORNATE NOUN (say "the purple cow"), then a PREPOSITION
(say “without"), and once again, we hit the recursion. So we hang onto our hats descend
one more level. To avoid complexity, let's assume that this the pathway we take is the
direct one just ORNATE NOUN. For example: we might get "horns". We hit the node
END in this call on FANCY NOUN which amounts to popping out, and so we go to our
stack to find the return address. It tells us that we were in the middle of executing
FANCY NOUN one level up-and so we resume there. This yields "the purple cow
without horns". On this level, too, we hit END, and so we pop up once more, this finding
ourselves in need of a VERB-so let's choose "gobbled". This ends highest-level call on
FANCY NOUN, with the result that the phrase

       "the strange bagels that the purple cow without horns gobbled"

will get passed upwards to the patient SENTENCE, as we pop for the last time.
        As you see, we didn't get into any infinite regress. The reason is tl least one
pathway inside the RTN FANCY NOUN does not involve recursive calls on FANCY
NOUN itself. Of course, we could have perversely insisted on always choosing the
bottom pathway inside FANCY NOUN then we would never have gotten finished, just
as the acronym "GOD” never got fully expanded. But if the pathways are chosen at
random, an infinite regress of that sort will not happen.

                      "Bottoming Out" and Heterarchies
        This is the crucial fact which distinguishes recursive definitions from circular
ones. There is always some part of the definition which avoids reference, so that the
action of constructing an object which satisfies the definition will eventually "bottom
out".
        Now there are more oblique ways of achieving recursivity in RTNs than by self-
calling. There is the analogue of Escher's Drawing (Fig. 135), where each of two
procedures calls the other, but not itself. For example, we could have an RTN named
CLAUSE, which calls FANCY NOUN whenever it needs an object for a transitive verb,
and conversely, the u path of FANCY NOUN could call RELATIVE PRONOUN and
then CLAUSE




Recursive Structures and Processes                                                    141

whenever it wants a relative clause. This is an example of indirect recursion. It is
reminiscent also of the two-step version of the Epimenides paradox.
         Needless to say, there can be a trio of procedures which call one another,
cyclically-and so on. There can be a whole family of RTN's which are all tangled up,
calling each other and themselves like crazy. A program which has such a structure in
which there is no single "highest level", or "monitor", is called a heterarchy (as
distinguished from a hierarchy). The term is due, I believe, to Warren McCulloch, one of
the first cyberneticists, and a reverent student of brains and minds.

                                 Expanding Nodes
One graphic way of thinking about RTN's is this. Whenever you are moving along some
pathway and you hit a node which calls on an RTN, you "expand" that node, which
means to replace it by a very small copy of the RTN it calls (see Fig. 28). Then you
proceed into the very small RTN,




      FIGURE 28. The FANCY NOUN RTN with one node recursively expanded

When you pop out of it, you are automatically in the right place in the big one. While in
the small one, you may wind up constructing even more miniature RTN's. But by
expanding nodes only when you come across them, you avoid the need to make an
infinite diagram, even when an RTN calls itself.
        Expanding a node is a little like replacing a letter in an acronym by the word it
stands for. The "GOD" acronym is recursive but has the defect-or advantage-that you
must repeatedly expand the `G'; thus it never bottoms out. When an RTN is implemented
as a real computer program, however, it always has at least one pathway which avoids
recursivity (direct or indirect) so that infinite regress is not created. Even the most
heterarchical program structure bottoms out-otherwise it couldn't run! It would just be
constantly expanding node after node, but never performing any action.




Recursive Structures and Processes                                                   142

                      Diagram G and Recursive Sequences
        Infinite geometrical structures can be defined in just this way-that is by expanding
node after node. For example, let us define an infinite diagram called "Diagram G". To
do so, we shall use an implicit representation. In two nodes, we shall write merely the
letter `G', which, however, will stand for an entire copy of Diagram G. In Figure 29a,
Diagram G is portrayed implicitly. Now if we wish to see Diagram G more explicitly, we
expand each of the two G's-that is, we replace them by the same diagram, only reduced
in scale (see Fig. 29b). This "second-order" version of Diagram gives us an inkling of
what the final, impossible-to-realize Diagram G really looks like. In Figure 30 is shown a
larger portion of Diagram G, where all the nodes have been numbered from the bottom
up, and from left to right. Two extra nodes-numbers -- 1 and 2--- have been inserted at
the bottom
        This infinite tree has some very curious mathematical properties Running up its
right-hand edge is the famous sequence of Fibonacci numbers.

       1,   1, 2, 3, 5, 8, 13, 21, 34, 55, 89,           144,    233,

discovered around the year 1202 by Leonardo of Pisa, son of Bonaccio, ergo "Filius
Bonacci", or "Fibonacci" for short. These numbers are best

FIGURE 29. (a) Diagram G, unexpanded.                (c) Diagram H, unexpanded
           (b) Diagram G, expanded once.             (d) Diagram H, expanded once




Recursive Structures and Processes                                                      143

             FIGURE 30. Diagram G, further expanded and with numbered nodes.

       defined recursively by the pair of formulas

       FIBO(n) = FIBO(n- 1) + FIBO(n-2)              for n > 2

              FIBO(l) = FIBO(2) = 1

Notice how new Fibonacci numbers are defined in terms of previous Fibonacci numbers.
We could represent this pair of formulas in an RTN (see Fig. 31).




                         FIGURE 31. An RTN for Fibonacci numbers.


Thus you can calculate FIBO(15) by a sequence of recursive calls on the procedure
defined by the RTN above. This recursive definition bottoms out when you hit FIBO(1)
or FIBO(2) (which are given explicitly) after you have worked your way backwards
through descending values of n. It is slightly awkward to work your way backwards,
when you could just as well work your way forwards, starting with FIBO(l) and FIBO(2)
and always adding the most recent two values, until you reach FIBO(15). That way you
don't need to keep track of a stack.
        Now Diagram G has some even more surprising properties than this. Its entire
structure can be coded up in a single recursive definition, as follows:


Recursive Structures and Processes                                               144

       G(n) = n - G(G(n- 1)) for n > 0

               G(O) = 0

How does this function G(n) code for the tree-structure? Quite simply you construct a
tree by placing G(n) below n, for all values of n, you recreate Diagram G. In fact, that is
how I discovered Diagram G in the place. I was investigating the function G, and in
trying to calculate its values quickly, I conceived of displaying the values I already knew
in a tree. T surprise, the tree turned out to have this extremely orderly recursive
geometrical description.
        What is more wonderful is that if you make the analogous tree function H(n)
defined with one more nesting than G—

       H(n) = n - H(H(H(n - 1)))      for n > 0

               H(0) = 0

--then the associated "Diagram H" is defined implicitly as shown in Figure 29c. The
right-hand trunk contains one more node; that is the difference. The first recursive
expansion of Diagram H is shown in Figure 29d. And so it goes, for any degree of
nesting. There is a beautiful regularity to the recursive geometrical structures, which
corresponds precisely to the recursive algebraic definitions.
        A problem for curious readers is: suppose you flip Diagram G around as if in a
mirror, and label the nodes of the new tree so they increase left to right. Can you find a
recursive algebraic definition for this "flip-tree. What about for the "flip" of the H-tree?
Etc.?
        Another pleasing problem involves a pair of recursively intertwined functions
F(n) and M(n) -- "married" functions, you might say -- defined this way:

       F(n) = n - M(F(n- 1))
                                                     For n > 0
       M(n) = n - F(M(n- 1))

       F(0) = 1, and M(0) = 0

       The RTN's for these two functions call each other and themselves as well. The
problem is simply to discover the recursive structures of Diagram F; and Diagram M.
They are quite elegant and simple.

                                 A Chaotic Sequence
       One last example of recursion in number theory leads to a small my Consider the
following recursive definition of a function:

       Q(n) = Q(n - Q(n- 1)) + Q(n - Q(n-2))         for n > 2
              Q(1) = Q(2) = 1.


Recursive Structures and Processes                                                      145

It is reminiscent of the Fibonacci definition in that each new value is a sum of two
previous values-but not of the immediately previous two values. Instead, the two
immediately previous values tell how far to count back to obtain the numbers to be added
to make the new value! The first 17 Q-numbers run as follows:

1, 1, 2, 3, 3, 4, 5, 5, 6, 6, 6, 8, 8, 8, 10, 9, 10, . . . .
 .. .                5 + 6 = 11               how far to move to the left
                      New term

To obtain the next one, move leftwards (from the three dots) respectively 10 and 9 terms;
you will hit a 5 and a 6, shown by the arrows. Their sum-1 l-yields the new value: Q(18).
This is the strange process by which the list of known Q-numbers is used to extend itself.
The resulting sequence is, to put it mildly, erratic. The further out you go, the less sense it
seems to make. This is one of those very peculiar cases where what seems to be a
somewhat natural definition leads to extremely puzzling behavior: chaos produced in a
very orderly manner. One is naturally led to wonder whether the apparent chaos conceals
some subtle regularity. Of course, by definition, there is regularity, but what is of interest
is whether there is another way of characterizing this sequence-and with luck, a
nonrecursive way.

                          Two Striking Recursive Graphs
The marvels of recursion in mathematics are innumerable, and it is not my purpose to
present them all. However, there are a couple of particularly striking examples from my
own experience which I feel are worth presenting. They are both graphs. One came up in
the course of some number-theoretical investigations. The other came up in the course of
my Ph.D. thesis work, in solid state physics. What is truly fascinating is that the graphs
are closely related.
         The first one (Fig. 32) is a graph of a function which I call INT(x). It is plotted
here for x between 0 and 1. For x between any other pair of integers n and n + 1, you just
find INT(x-n), then add n back. The structure of the plot is quite jumpy, as you can see. It
consists of an infinite number of curved pieces, which get smaller and smaller towards
the corners-and incidentally, less and less curved. Now if you look closely at each such
piece, you will find that it is actually a copy of the full graph, merely curved! The
implications are wild. One of them is that the graph of INT consists of nothing but copies
of itself, nested down infinitely deeply. If you pick up any piece of the graph, no matter
how small, you are holding a complete copy of the whole graph-in fact, infinitely many
copies of it!

The fact that INT consists of nothing but copies of itself might make you think it is too
ephemeral to exist. Its definition sounds too circular.




Recursive Structures and Processes                                                         146

FIGURE 32. Graph of the function INT(x). There is a jump discontinuity at every rat
value of x.

How does it ever get off the ground? That is a very interesting matter. main thing to
notice is that, to describe INT to someone who hasn't see it will not suffice merely to say,
"It consists of copies of itself." The o half of the story-the nonrecursive half-tells where
those copies lie in the square, and how they have been deformed, relative to the full
graph. Only the combination of these two aspects of INT will specify structure of INT. It
is exactly as in the definition of Fibonacci number where you need two lines-one to
define the recursion, the other to de the bottom (i.e., the values at the beginning). To be
very concrete, if make one of the bottom values 3 instead of 1, you will produce a
completely different sequence, known as the Lucas sequence:

       1, 3, 4 , 7,     11,   18,   29, 47, 76, 123, .. .
       the "bottom"                 29 + 47 = 76
                                     same recursive rule
                                     as for the Fibonacci numbers




Recursive Structures and Processes                                                      147

         What corresponds to the bottom in the definition of INT is a picture (Fig. 33a)
composed of many boxes, showing where the copies go, and how they are distorted. I call
it the "skeleton" of INT. To construct INT from its skeleton, you do the following. First,
for each box of the skeleton, you do two operations: (1) put a small curved copy of the
skeleton inside the box, using the curved line inside it as a guide; (2) erase the containing
box and its curved line. Once this has been done for each box of the original skeleton,
you are left with many "baby" skeletons in place of one big one. Next you repeat the
process one level down, with all the baby skeletons. Then again, again, and again ... What
you approach in the limit is an exact graph of INT, though you never get there. By
nesting the skeleton inside itself over and over again, you gradually construct the graph
of INT "from out of nothing". But in fact the "nothing" was not nothing-it was a picture.
         To see this even more dramatically, imagine keeping the recursive part of the
definition of INT, but changing the initial picture, the skeleton. A variant skeleton is
shown in Figure 33b, again with boxes which get smaller and smaller as they trail off to
the four corners. If you nest this second skeleton inside itself over and over again, you
will create the key graph from my Ph.D. thesis, which I call Gplot (Fig. 34). (In fact,
some complicated distortion of each copy is needed as well-but nesting is the basic idea.).
         Gplot is thus a member of the INT-family. It is a distant relative, because its
skeleton is quite different from-and considerably more complex than-that of INT.
However, the recursive part of the definition is identical, and therein lies the family tie.
I should not keep you too much in the dark about the origin of these beautiful graphs.
INT-standing for "interchange"-comes from a problem involving "Eta-sequences", which
are related to continued fractions. The basic idea behind INT is that plus and minus signs
are interchanged in a certain kind of continued fraction. As a consequence, INT(INT(x))
= x. INT has the property that if x is rational, so is INT(x); if x is quadratic, so is INT(x).
I do not know if this trend holds for higher algebraic degrees. Another lovely feature of
INT is that at all rational values of x, it has a jump discontinuity, but at all irrational
values of x, it is continuous.
         Gplot comes from a highly idealized version of the question, "What are the
allowed energies of electrons in a crystal in a magnetic field?" This problem is interesting
because it is a cross between two very simple and fundamental physical situations: an
electron in a perfect crystal, and an electron in a homogeneous magnetic field. These two
simpler problems are both well understood, and their characteristic solutions seem almost
incompatible with each other. Therefore, it is of quite some interest to see how nature
manages to reconcile the two. As it happens, the crystal without-magnetic-field situation
and the magnetic-field-without-crystal situation do have one feature in common: in each
of them, the electron behaves periodically in time. It turns out that when the two
situations are combined, the ratio of their two time periods is the key parameter. In fact,
that ratio holds all the information about the distribution of allowed electron energies-but
it only gives up its secret upon being expanded into a continued fraction.




Recursive Structures and Processes                                                         148

Recursive Structures and Processes   149

        Gplot shows that distribution. The horizontal axis represents energy, and the
vertical axis represents the above-mentioned ratio of time periods, which we can call "α".
At the bottom, a is zero, and at the top a is unity. When a is zero, there is no magnetic
field. Each of the line segments making up Gplot is an "energy band"-that is, it represents
allowed values of energy. The empty swaths traversing Gplot on all different size scales
are therefore regions of forbidden energy. One of the most startling properties of Gplot is
that when a is rational (say p/q in lowest terms), there are exactly q such bands (though
when q is even, two of them "kiss" in the middle). And when a is irrational, the bands
shrink to points, of which there are infinitely many, very sparsely distributed in a so-
called "Cantor set" -- another recursively defined entity which springs up in topology.
        You might well wonder whether such an intricate structure would ever show up in
an experiment. Frankly, I would be the most surprised person in the world if Gplot came
out of any experiment. The physicality of Gplot lies in the fact that it points the way to
the proper mathematical treatment of less idealized problems of this sort. In other words,
Gplot is purely a contribution to theoretical physics, not a hint to experimentalists as to
what to expect to see! An agnostic friend of mine once was so struck by Gplot's infinitely
many infinities that he called it "a picture of God", which I don't think is blasphemous at
all.

                    Recursion at the Lowest Level of Matter
        We have seen recursion in the grammars of languages, we have seen recursive
geometrical trees which grow upwards forever, and we have seen one way in which
recursion enters the theory of solid state physics. Now we are going to see yet another
way in which the whole world is built out of recursion. This has to do with the structure
of elementary particles: electrons, protons, neutrons, and the tiny quanta of
electromagnetic radiation called "photons". We are going to see that particles are-in a
certain sense which can only be defined rigorously in relativistic quantum mechanics --
nested inside each other in a way which can be described recursively, perhaps even by
some sort of "grammar".
        We begin with the observation that if particles didn't interact with each other,
things would be incredibly simple. Physicists would like such a world because then they
could calculate the behavior of all particles easily (if physicists in such a world existed,
which is a doubtful proposition). Particles without interactions are called bare particles,
and they are purely hypothetical creations; they don't exist.
        Now when you "turn on" the interactions, then particles get tangled up together in
the way that functions F and M are tangled together, or married people are tangled
together. These real particles are said to be renormalized-an ugly but intriguing term.
What happens is that no particle can even be defined without referring to all other
particles, whose definitions in turn depend on the first particles, etc. Round and round, in
a never-ending loop.




Recursive Structures and Processes                                                      150

Figure 34. Gplot; a recursive graph, showing energy bands for electrons in an idealized
crystal in a magnetic field, α representing magnetic field strength, runs vertically from 0
to 1. Energy runs horizontally. The horizontal line segments are bands of allowed
electron energies.


Recursive Structures and Processes                                                     151

        Let us be a little more concrete, now. Let's limit ourselves to only two kinds of
particles: electrons and photons. We'll also have to throw in the electron's antiparticle, the
positron. (Photons are their own antiparticles.) Imagine first a dull world where a bare
electron wishes to propagate from point A to point B, as Zeno did in my Three-Part
Invention. A physicist would draw a picture like this:




There is a mathematical expression which corresponds to this line and its endpoints, and
it is easy to write down. With it, a physicist can understand the behavior of the bare
electron in this trajectory.
        Now let us "turn on" the electromagnetic interaction, whereby electrons and
photons interact. Although there are no photons in the scene, there will nevertheless be
profound consequences even for this simple trajectory. In particular, our electron now
becomes capable of emitting and then reabsorbing virtual photons-photons which flicker
in and out of existence before they can be seen. Let us show one such process:




Now as our electron propagates, it may emit and reabsorb one photon after another, or it
may even nest them, as shown below:




The mathematical expressions corresponding to these diagrams-called "Feynman
diagrams"-are easy to write down, but they are harder to calculate than that for the bare
electron. But what really complicates matters is that a photon (real or virtual) can decay
for a brief moment into an electron-positron pair. Then these two annihilate each other,
and, as if by magic, the original photon reappears. This sort of process is shown below:




The electron has a right-pointing arrow, while the positron's arrow points leftwards.




Recursive Structures and Processes                                                        152

        As you might have anticipated, these virtual processes can be inside each other to
arbitrary depth. This can give rise to some complicated-looking drawings, such as the one
in Figure 35. In that man diagram, a single electron enters on the left at A, does some an
acrobatics, and then a single electron emerges on the right at B. outsider who can't see the
inner mess, it looks as if one electron peacefully sailed from A to B. In the diagram, you
can see how el lines can get arbitrarily embellished, and so can the photon lines diagram
would be ferociously hard to calculate.
        .




FIGURE 35. A Feynman diagram showing the propagation of a renormalized electron
from A to B. In this diagram, time increases to the right. Therefore, in the segments
where the electron’s arrow points leftwards, it is moving "backwards in time". A more
intuitive way to say this is that an antielectron (positron) is moving forwards in time.
Photons are their own antiparticles; hence their lines have no need of arrows.

        There is a sort of "grammar" to these diagrams, that only certain pictures to be
realized in nature. For instance, the one be impossible:




You might say it is not a "well-formed" Feynman diagram. The gram a result of basic
laws of physics, such as conservation of energy, conservation of electric charge, and so
on. And, like the grammars of l - languages, this grammar has a recursive structure, in
that it allow' nestings of structures inside each other. It would be possible to drat set of
recursive transition networks defining the "grammar" of the electromagnetic interaction.
         When bare electrons and bare photons are allowed to interact ii arbitrarily tangled
ways, the result is renormalized electrons and ph Thus, to understand how a real, physical
electron propagates from A to B,




Recursive Structures and Processes                                                      153

the physicist has to be able to take a sort of average of all the infinitely many different
possible drawings which involve virtual particles. This is Zeno with a vengeance!
        Thus the point is that a physical particle-a renormalized particle involves (1) a
bare particle and (2) a huge tangle of virtual particles, inextricably wound together in a
recursive mess. Every real particle's existence therefore involves the existence of
infinitely many other particles, contained in a virtual "cloud" which surrounds it as it
propagates. And each of the virtual particles in the cloud, of course, also drags along its
own virtual cloud, and so on ad infinitum.
        Particle physicists have found that this complexity is too much to handle, and in
order to understand the behavior of electrons and photons, they use approximations
which neglect all but fairly simple Feynman diagrams. Fortunately, the more complex a
diagram, the less important its contribution. There is no known way of summing up all of
the infinitely many possible diagrams, to get an expression for the behavior of a fully
renormalized, physical electron. But by considering roughly the simplest hundred
diagrams for certain processes, physicists have been able to predict one value (the so-
called g-factor of the muon) to nine decimal places -- correctly!
        Renormalization takes place not only among electrons and photons. Whenever
any types of particle interact together, physicists use the ideas of renormalization to
understand the phenomena. Thus protons and neutrons, neutrinos, pi-mesons, quarks-all
the beasts in the subnuclear zoo they all have bare and renormalized versions in physical
theories. And from billions of these bubbles within bubbles are all the beasts and baubles
of the world composed.

                                 Copies and Sameness
         Let us now consider Gplot once again. You will remember that in the
Introduction, we spoke of different varieties of canons. Each type of canon exploited
some manner of taking an original theme and copying it by an isomorphism, or
information-preserving transformation. Sometimes the copies were upside down,
sometimes backwards, sometimes shrunken or expanded ... In Gplot we have all those
types of transformation, and more. The mappings between the full Gplot and the "copies"
of itself inside itself involve size changes, skewings, reflections, and more. And yet there
remains a sort of skeletal identity, which the eye can pick up with a bit of effort,
particularly after it has practiced with INT.
         Escher took the idea of an object's parts being copies of the object itself and made
it into a print: his woodcut Fishes and Scales (Fig. 36). Of course these fishes and scales
are the same only when seen on a sufficiently abstract plane. Now everyone knows that a
fish's scales aren't really small copies of the fish; and a fish's cells aren't small copies of
the fish; however, a fish's DNA, sitting inside each and every one of the fish's cells, is a
very convo-




Recursive Structures and Processes                                                         154

            FIGURE 36. Fish and Scales, by M. C. Escher (woodcut, 1959).

luted "copy" of the entire fish-and so there is more than a grain of truth to the Escher
picture.
        What is there that is the "same" about all butterflies? The mapping from one
butterfly to another does not map cell onto cell; rather, it m; functional part onto
functional part, and this may be partially on a macroscopic scale, partially on a
microscopic scale. The exact proportions of pa are not preserved; just the functional
relationships between parts. This is the type of isomorphism which links all butterflies in
Escher's wood engraving Butterflies (Fig. 37) to each other. The same goes for the more
abstract butterflies of Gplot, which are all linked to each other by mathematical mappings
that carry functional part onto functional part, but totally ignore exact line proportions,
angles, and so on.
        Taking this exploration of sameness to a yet higher plane of abstraction, we might
well ask, "What is there that is the `same' about all Esc l drawings?" It would be quite
ludicrous to attempt to map them piece by piece onto each other. The amazing thing is
that even a tiny section of an




Recursive Structures and Processes                                                     155

           FIGURE 37. Butterflies, by M. C. Escher (wood-engraving, 1950).

Escher drawing or a Bach piece gives it away. Just as a fish's DNA is contained inside
every tiny bit of the fish, so a creator's "signature" is contained inside every tiny section
of his creations. We don't know what to call it but "style" -- a vague and elusive word.
        We keep on running up against "sameness-in-differentness", and the question

                       When are two things the same?

It will recur over and over again in this book. We shall come at it from all sorts of skew
angles, and in the end, we shall see how deeply this simple question is connected with the
nature of intelligence.
         That this issue arose in the Chapter on recursion is no accident, for recursion is a
domain where "sameness-in-differentness" plays a central role. Recursion is based on the
"same" thing happening on several differ-




Recursive Structures and Processes                                                       156

ent levels at once. But the events on different levels aren't exactly same-rather, we find
some invariant feature in them, despite many s in which they differ. For example, in the
Little Harmonic Labyrinth, all stories on different levels are quite unrelated-their
"sameness" reside only two facts: (1) they are stories, and (2) they involve the Tortoise
and Achilles. Other than that, they are radically different from each other.

      Programming and Recursion: Modularity, Loops, Procedures
One of the essential skills in computer programming is to perceive wl two processes are
the same in this extended sense, for that leads modularization-the breaking-up of a task
into natural subtasks. For stance, one might want a sequence of many similar operations
to be cart out one after another. Instead of writing them all out, one can write a h which
tells the computer to perform a fixed set of operations and then loop back and perform
them again, over and over, until some condition is satisfied. Now the body of the loop-the
fixed set of instructions to repeated-need not actually be completely fixed. It may vary in
so predictable way.
         An example is the most simple-minded test for the primality o natural number N,
in which you begin by trying to divide N by 2, then 3, 4, 5, etc. until N - 1. If N has
survived all these tests without be divisible, it's prime. Notice that each step in the loop is
similar to, but i the same as, each other step. Notice also that the number of steps varies
with N-hence a loop of fixed length could never work as a general test primality. There
are two criteria for "aborting" the loop: (1) if so number divides N exactly, quit with
answer "NO"; (2) if N - 1 is react as a test divisor and N survives, quit with answer
"YES".
         The general idea of loops, then, is this: perform some series of related steps over
and over, and abort the process when specific conditions are n Now sometimes, the
maximum number of steps in a loop will be known advance; other times, you just begin,
and wait until it is aborted. The second type of loop -- which I call a free loop -- is
dangerous, because criterion for abortion may never occur, leaving the computer in a so-
cal "infinite loop". This distinction between bounded loops and free loops is one the most
important concepts in all of computer science, and we shall dev an entire Chapter to it:
"BlooP and FlooP and G1ooP".
         Now loops may be nested inside each other. For instance, suppose t we wish to
test all the numbers between 1 and 5000 for primality. We c write a second loop which
uses the above-described test over and over starting with N = I and finishing with N =
5000. So our program i have a "loop-the-loop" structure. Such program structures are
typical – in fact they are deemed to be good programming style. This kind of nest loop
also occurs in assembly instructions for commonplace items, and such activities as
knitting or crocheting-in which very small loops are




Recursive Structures and Processes                                                         157

repeated several times in larger loops, which in turn are carried out repeatedly ... While
the result of a low-level loop might be no more than couple of stitches, the result of a
high-level loop might be a substantial portion of a piece of clothing.
        In music, too, nested loops often occur-as, for instance, when a scale (a small
loop) is played several times in a row, perhaps displaced in pitch each new time. For
example, the last movements of both the Prokofiev fifth piano concerto and the
Rachmaninoff second symphony contain extended passages in which fast, medium, and
slow scale-loops are played simultaneously by different groups of instruments, to great
effect. The Prokofiev scales go up; the Rachmaninoff-scales, down. Take your pick.
        A more general notion than loop is that of subroutine, or procedure, which we
have already discussed somewhat. The basic idea here is that a group of operations are
lumped together and considered a single unit with a name-such as the procedure
ORNATE NOUN. As we saw in RTN's, procedures can call each other by name, and
thereby express very concisely sequences of operations which are to be carried out. This
is the essence of modularity in programming. Modularity exists, of course, in hi-fi
systems, furniture, living cells, human society-wherever there is hierarchical
organization.
        More often than not, one wants a procedure which will act variably, according to
context. Such a procedure can either be given a way of peering out at what is stored in
memory and selecting its actions accordingly, or it can be explicitly fed a list of
parameters which guide its choice of what actions to take. Sometimes both of these
methods are used. In RTN terminology, choosing the sequence of actions to carry out
amounts to choosing which pathway to follow. An RTN which has been souped up with
parameters and conditions that control the choice of pathways inside it is called an
Augmented Transition Network (ATN). A place where you might prefer ATN's to RTN's
is in producing sensible-as distinguished from nonsensical-English sentences out of raw
words, according to a grammar represented in a set of ATN's. The parameters and
conditions would allow you to insert various semantic constraints, so that random
juxtapositions like "a thankless brunch" would be prohibited. More on this in Chapter
XVIII, however.


                          Recursion in Chess Programs
A classic example of a recursive procedure with parameters is one for choosing the "best"
move in chess. The best move would seem to be the one which leaves your opponent in
the toughest situation. Therefore, a test for goodness of a move is simply this: pretend
you've made the move, and now evaluate the board from the point of view of your
opponent. But how does your opponent evaluate the position? Well, he looks for his best
move. That is, he mentally runs through all possible moves and evaluates them from what
he thinks is your point of view, hoping they will look bad to you. But




Recursive Structures and Processes                                                    158

notice that we have now defined "best move" recursively, simply maxim that what is best
for one side is worst for the other. The procedure which looks for the best move operates
by trying a move and then calling on itself in the role of opponent! As such, it tries
another n calls on itself in the role of its opponent's opponent-that is, its
         This recursion can go several levels deep-but it's got to bottom out somewhere!
How do you evaluate a board position without looking There are a number of useful
criteria for this purpose, such as si number of pieces on each side, the number and type of
pieces undo the control of the center, and so on. By using this kind of evaluation at the
bottom, the recursive move-generator can pop back upwards an( evaluation at the top
level of each different move. One of the parameters in the self-calling, then, must tell
how many moves to look ahead. TI most call on the procedure will use some externally
set value parameter. Thereafter, each time the procedure recursively calls must decrease
this look-ahead parameter by 1. That way, w parameter reaches zero, the procedure will
follow the alternate pathway -- the non-recursive evaluation.
         In this kind of game-playing program, each move investigate the generation of a
so-called "look-ahead tree", with the move trunk, responses as main branches, counter-
responses as subsidiary branches, and so on. In Figure 38 I have shown a simple look-
ahead tree depicting the start of a tic-tar-toe game. There is an art to figuring to avoid
exploring every branch of a look-ahead tree out to its tip. trees, people-not computers-
seem to excel at this art; it is known that top-level players look ahead relatively little,
compared to most chess programs – yet the people are far better! In the early days of
compute people used to estimate that it would be ten years until a computer (or

FIGURE 38. The branching tree of moves and countermoves at the start of c tic-tac-toe.




Recursive Structures and Processes                                                     159

 program) was world champion. But after ten years had passed, it seemed that the day a
 computer would become world champion was still more than ten years away ... This is
               just one more piece of evidence for the rather recursive

      Hofstadter's Law: It always takes longer than you expect, even when you take into
             account Hofstadter's Law.

                          Recursion and Unpredictability
Now what is the connection between the recursive processes of this Chapter, and the
recursive sets of the preceding Chapter? The answer involves the notion of a recursively
enumerable set. For a set to be r.e. means that it can be generated from a set of starting
points (axioms), by the repeated application of rules of inference. Thus, the set grows and
grows, each new element being compounded somehow out of previous elements, in a sort
of "mathematical snowball". But this is the essence of recursion-something being defined
in terms of simpler versions of itself, instead of explicitly. The Fibonacci numbers and
the Lucas numbers are perfect examples of r.e. sets-snowballing from two elements by a
recursive rule into infinite sets. It is just a matter of convention to call an r.e. set whose
complement is also r.e. "recursive".
        Recursive enumeration is a process in which new things emerge from old things
by fixed rules. There seem to be many surprises in such processes-for example the
unpredictability of the Q-sequence. It might seem that recursively defined sequences of
that type possess some sort of inherently increasing complexity of behavior, so that the
further out you go, the less predictable they get. This kind of thought carried a little
further suggests that suitably complicated recursive systems might be strong enough to
break out of any predetermined patterns. And isn't this one of the defining properties of
intelligence? Instead of just considering programs composed of procedures which can
recursively call themselves, why not get really sophisticated, and invent programs which
can modify themselves-programs which can act on programs, extending them, improving
them, generalizing them, fixing them, and so on? This kind of "tangled recursion"
probably lies at the heart of intelligence.




Recursive Structures and Processes                                                        160

                                       Canon
                            by Intervallic Augmentation

             Achilles and the Tortoise have just finished a delicious Chinese banquet for
             two, at the best Chinese restaurant in town.


Achilles: You wield a mean chopstick, Mr. T.
Tortoise: I ought to. Ever since my youth, I have had a fondness for Oriental cuisine. And you-
       did you enjoy your meal, Achilles? Achilles: Immensely. I'd not eaten Chinese food
       before. This meal was a splendid introduction. And now, are you in a hurry to go, or shall
       we just sit here and talk a little while?
Tortoise: I'd love to talk while we drink our tea. Waiter!

       (A waiter comes up.)

       Could we have our bill, please, and some more tea?

       (The waiter rushes off.)

Achilles: You may know more about Chinese cuisine than I do, Mr.T, I'll bet I know more about
       Japanese poetry than you do. Have you ever read any haiku?
Tortoise: I'm afraid not. What is a haiku?
Achilles: A haiku is a Japanese seventeen-syllable poem-or minipoem rather, which is evocative
       in the same way, perhaps, as a fragrant petal is, or a lily pond in a light drizzle. It
       generally consists of groups of: of five, then seven, then five syllables.
Tortoise: Such compressed poems with seventeen syllables can't much meaning ...
Achilles: Meaning lies as much in the mind of the reader as i haiku.
Tortoise: Hmm ... That's an evocative statement.

        (The waiter arrives with their bill, another pot of tea, and two fortune cookies.)

       Thank you, waiter. Care for more tea, Achilles?
Achilles: Please. Those little cookies look delicious. (Picks one up, bites I into it and begins to
       chew.) Hey! What's this funny thing inside? A piece of paper?
Tortoise: That's your fortune, Achilles. Many Chinese restaurants give out fortune cookies with
       their bills, as a way of softening the blow. I frequent Chinese restaurants, you come to
       think of fortune cookies

        less as cookies than as message bearers Unfortunately you seem to have swallowed some
       of your fortune. What does the rest say?
Achilles: It's a little strange, for all the letters are run together, with no spaces in between.
       Perhaps it needs decoding in some way? Oh, now I see. If you put the spaces back in
       where they belong, it says, "ONE WAR TWO EAR EWE". I can't quite make head or tail
       of that. Maybe it was a haiku-like poem, of which I ate the majority of syllables.
Tortoise: In that case, your fortune is now a mere 5/17-haiku. And a curious image it evokes. If
       5/17-haiku is a new art form, then I'd say woe, 0, woe are we ... May I look at it?
Achilles (handing the Tortoise the small slip of paper): Certainly.
Tortoise: Why, when I "decode" it, Achilles, it comes out completely different! It's not a 5/17-
       haiku at all. It is a six-syllable message which says, "0 NEW ART WOE ARE WE". That
       sounds like an insightful commentary on the new art form of 5/17-haiku.
Achilles: You're right. Isn't it astonishing that the poem contains its own commentary!
Tortoise: All I did was to shift the reading frame by one unit-that is, shift all the spaces one unit
       to the right.
Achilles: Let's see what your fortune says, Mr. Tortoise.
Tortoise (deftly splitting open his cookie, reads): "Fortune lies as much in the hand of the eater as
       in the cookie."
Achilles: Your fortune is also a haiku, Mr. Tortoise-at least it's got seventeen syllables in the 5-7-
       5 form.
Tortoise: Glory be! I would never have noticed that, Achilles. It's the kind of thing only you
       would have noticed. What struck me more is what it says-which, of course, is open to
       interpretation.
Achilles: I guess it just shows that each of us has his own characteristic way of interpreting
       messages which we run across ...

       (Idly, Achilles gazes at the tea leaves on the bottom of his empty teacup.)

Tortoise: More tea, Achilles?
Achilles: Yes, thank you. By the way, how is your friend the Crab? I have been thinking about
       him a lot since you told me of your peculiar phonograph-battle.
Tortoise: I have told him about you, too, and he is quite eager to meet you. He is getting along
       just fine. In fact, he recently made a new acquisition in the record player line: a rare type
       of jukebox.
Achilles: Oh, would you tell me about it? I find jukeboxes, with their flashing colored lights and
       silly songs, so quaint and reminiscent of bygone eras.
Tortoise: This jukebox is too large to fit in his house, so he had a shed specially built in back for
       it.
Achilles: I can't imagine why it would be so large, unless it has an unusually large selection of
       records. Is that it?
Tortoise: As a matter of fact, it has exactly one record.

Achilles: What? A jukebox with only one record? That's a contradiction in terms. Why is the
        jukebox so big, then? Is its single record gigantic -- twenty feet in diameter?
Tortoise: No, it's just a regular jukebox-style record.
Achilles: Now, Mr. Tortoise, you must be joshing me. After all, what I
of a jukebox is it that has only a single song?
Tortoise: Who said anything about a single song, Achilles?
Achilles: Every- jukebox I've ever run into obeyed the fundamental jukebox-axiom: "One record,
        one song".
Tortoise: This jukebox is different, Achilles. The one record sits vertically, suspended, and
        behind it there is a small but elaborate network of overhead rails, from which hang
        various record players. When push a pair of buttons, such as B-1, that selects one of the
        record players. This triggers an automatic mechanism that starts the record player
        squeakily rolling along the rusty tracks. It gets shunted alongside the record-then it clicks
        into playing position.
Achilles: And then the record begins spinning and music comes out -- right?
Tortoise: Not quite. The record stands still-it's the record player which rotates.
Achilles: I might have known. But how, if you have but one record to play can you get more than
        one song out of this crazy contraption?
Tortoise: I myself asked the Crab that question. He merely suggested I try it out. So I fished a
        quarter from my pocket (you get three plays for a quarter), stuffed it in the slot, and hit
        buttons B-1, then C-3 then B-10-all just at random.
Achilles: So phonograph B-1 came sliding down the rail, I suppose, plugged itself into the
        vertical record, and began spinning?
Tortoise: Exactly. The music that came out was quite agreeable, based the famous old tune B-A-
        C-H, which I believe you remember.




Achilles: Could I ever forget it?
Tortoise: This was record player B-1. Then it finished, and was s rolled back into its hanging
       position, so that C-3 could be slid into position.
Achilles: Now don't tell me that C-3 played another song?
Tortoise: It did just that.
Achilles: Ah, I understand. It played the flip side of the first song, or another band on the same
       side.
Tortoise: No, the record has grooves only on one side, and has only a single band.

Achilles: I don't understand that at all. You CAN'T pull different songs out of the same record!
Tortoise: That's what I thought until I saw Mr. Crab's jukebox. Achilles: How did the second
       song go?
Tortoise: That's the interesting thing ... It was a song based on the melody C-A-G-E.
Achilles: That's a totally different melody!
Tortoise: True.
Achilles: And isn't John Cage a composer of modern music? I seem to remember reading about
       him in one of my books on haiku.
Tortoise: Exactly. He has composed many celebrated pieces, such as 4'33", a three-movement
       piece consisting of silences of different lengths. It's wonderfully expressive-if you like
       that sort of thing.
Achilles: I can see where if I were in a loud and brash cafe I might gladly pay to hear Cage's
       4'33" on a jukebox. It might afford some relief!
Tortoise: Right-who wants to hear the racket of clinking dishes and jangling silverware? By the
       way, another place where 4'33" would come in handy is the Hall of Big Cats, at feeding
       time.
Achilles: Are you suggesting that Cage belongs in the zoo? Well, I guess that makes some sense.
       But about the Crab's jukebox ... I am baffled. How could both "BACH" and "CAGE" be
       coded inside a single record at once?
Tortoise: You may notice that there is some relation between the two, Achilles, if you inspect
       them carefully. Let me point the way. What do you get if you list the successive intervals
       in the melody B-A-C-H?
Achilles: Let me see. First it goes down one semitone, from B to A (where B is taken the
       German way); then it rises three semitones to C; and finally it falls one semitone, to H.
       That yields the pattern:

                                           -1, +3, -1.

Tortoise: Precisely. What about C-A-G-E, now?
Achilles: Well, in this case, it begins by falling three semitones, then ten semitones (nearly an
       octave), and finally falls three more semitones. That means the pattern is:

                                          -3, +10, -3.

       It's very much like the other one, isn't it?
Tortoise: Indeed it is. They have exactly the same "skeleton", in a certain sense. You can make
       C-A-G-E out of B-A-C-H by multiplying all the intervals by 31/3, and taking the nearest
       whole number.
Achilles: Well, blow me down and pick me up! So does that mean that only

       some sort of skeletal code is present in the grooves, and that the various record players
       add their own interpretations to that code?
Tortoise: I don't know, for sure. The cagey Crab wouldn't fill me in on the details. But I did get
       to hear a third song, when record player B-1 swiveled into place.
Achilles: How did it go?
Tortoise: The melody consisted of enormously wide intervals, and we B-C-A-H.




       The interval pattern in semitones was:

                              -10,    +33,      -10.

       It can be gotten from the CAGE pattern by yet another multiplication by 3\%3, and
       rounding to whole numbers.
Achilles: Is there a name for this kind of interval multiplication?
Tortoise: One could call it "intervallic augmentation". It is similar to tl canonic device of
       temporal augmentation, where all the time values notes in a melody get multiplied by
       some constant. There, the effect just to slow the melody down. Here, the effect is to
       expand the melodic range in a curious way.
Achilles: Amazing. So all three melodies you tried were intervallic augmentations of one single
       underlying groove-pattern in the record:
Tortoise: That's what I concluded.
Achilles: I find it curious that when you augment BACH you get CAGE and when you augment
       CAGE over again, you get BACH back, except jumbled up inside, as if BACH had an
       upset stomach after passing through the intermediate stage of CAGE.
Tortoise: That sounds like an insightful commentary on the new art form of Cage.


                       The Location of Meaning
                   When Is One Thing Not Always the Same?
LAST CHAPTER, WE came upon the question, "When are two things the same?" In this
Chapter, we will deal with the flip side of that question: "When is one thing not always
the same?" The issue we are broaching is whether meaning can be said to be inherent in a
message, or whether meaning is always manufactured by the interaction of a mind or a
mechanism with a message-as in the preceding Dialogue. In the latter case, meaning
could not said to be located in any single place, nor could it be said that a message has
any universal, or objective, meaning, since each observer could bring its own meaning to
each message. But in the former case, meaning would have both location and
universality. In this Chapter, I want to present the case for the universality of at least
some messages, without, to be sure, claiming it for all messages. The idea of an
"objective meaning" of a message will turn out to be related, in an interesting way, to the
simplicity with which intelligence can be described.

Information-Bearers and Information- Revealers
I'll begin with my favorite example: the relationship between records, music, and record
players. We feel quite comfortable with the idea that a record contains the same
information as a piece of music, because of the existence of record players, which can
"read" records and convert the groove-patterns into sounds. In other words, there is an
isomorphism between groove-patterns and sounds, and the record player is a mechanism
which physically realizes that isomorphism. It is natural, then, to think of the record as an
information-bearer, and the record-player as an information-revealer. A second example
of these notions is given by the pq-system. There, the "information-bearers" are the
theorems, and the "information-revealer" is the interpretation, which is so transparent that
we don't need any electrical machine to help us extract the information from pq-
theorems.
         One gets the impression from these two examples that isomorphisms and
decoding mechanisms (i.e., information-revealers) simply reveal information which is
intrinsically inside the structures, waiting to be "pulled out". This leads to the idea that
for each structure, there are certain pieces of information which can be pulled out of it,
while there are other pieces of information which cannot be pulled out of it. But what
does this phrase




The Location of Meaning                                                                  166

"pull out" really mean? How hard are you allowed to pull? There are c where by
investing sufficient effort, you can pull very recondite piece of information out of certain
structures. In fact, the pulling-out may inv such complicated operations that it makes you
feel you are putting in n information than you are pulling out.

                              Genotype and Phenotype
Take the case of the genetic information commonly said to reside in double helix of
deoxyribonucleic acid (DNA). A molecule of DNA – a genotype-is converted into a
physical organism-a phenotype-by a complex process, involving the manufacture of
proteins, the replication the DNA, the replication of cells, the gradual differentiation of
cell types and so on. Incidentally, this unrolling of phenotype from genotype epigenesis-
is the most tangled of tangled recursions, and in Chapter we shall devote our full attention
to it. Epigenesis is guided by a se enormously complex cycles of chemical reactions and
feedback loops the time the full organism has been constructed, there is not even remotest
similarity between its physical characteristics and its genotype.
        And yet, it is standard practice to attribute the physical structure of organism to
the structure of its DNA, and to that alone. The first evidence for this point of view came
from experiments conducted by Oswald A, in 1946, and overwhelming corroborative
evidence has since been amassed Avery's experiments showed that, of all the biological
molecules, only E transmits hereditary properties. One can modify other molecules it
organism, such as proteins, but such modifications will not be transmitted to later
generations. However, when DNA is modified, all successive generations inherit the
modified DNA. Such experiments show that the only of changing the instructions for
building a new organism is to change DNA-and this, in turn, implies that those
instructions must be cc somehow in the structure of the DNA.

                        Exotic and Prosaic Isomorphisms
Therefore one seems forced into accepting the idea that the DNA's structure contains the
information of the phenotype's structure, which is to the two are isomorphic. However,
the isomorphism is an exotic one, by w] I mean that it is highly nontrivial to divide the
phenotype and genotype into "parts" which can be mapped onto each other. Prosaic
isomorphic by contrast, would be ones in which the parts of one structure are easily
mappable onto the parts of the other. An example is the isomorphism between a record
and a piece of music, where one knows that to any so in the piece there exists an exact
"image" in the patterns etched into grooves, and one could pinpoint it arbitrarily
accurately, if the need arose Another prosaic isomorphism is that between Gplot and any
of its internal butterflies.




The Location of Meaning                                                                 167

The isomorphism between DNA structure and phenotype structure is anything but
prosaic, and the mechanism which carries it out physically is awesomely complicated.
For instance, if you wanted to find some piece of your DNA which accounts for the shape
of your nose or the shape of your fingerprint, you would have a very hard time. It would
be a little like trying to pin down the note in a piece of music which is the carrier of the
emotional meaning of the piece. Of course there is no such note, because the emotional
meaning is carried on a very high level, by large "chunks" of the piece, not by single
notes. Incidentally, such "chunks" are not necessarily sets of contiguous notes; there may
be disconnected sections which, taken together, carry some emotional meaning.
         Similarly, "genetic meaning"-that is, information about phenotype structure-is
spread all through the small parts of a molecule of DNA, although nobody understands
the language yet. (Warning: Understanding this "language" would not at all be the same
as cracking the Genetic Code, something which took place in the early 1960's. The
Genetic Code tells how to translate short portions of DNA into various amino acids.
Thus, cracking the Genetic Code is comparable to figuring out the phonetic values of the
letters of a foreign alphabet, without figuring out the grammar of the language or the
meanings of any of its words. The cracking of the Genetic Code was a vital step on the
way to extracting the meaning of DNA strands, but it was only the first on a long path
which is yet to be trodden.)


                               Jukeboxes and Triggers
The genetic meaning contained in DNA is one of the best possible examples of implicit
meaning. In order to convert genotype into phenotype, a set of mechanisms far more
complex than the genotype must operate on the genotype. The various parts of the
genotype serve as triggers for those mechanisms. A jukebox-the ordinary type, not the
Crab type!-provides a useful analogy here: a pair of buttons specifies a very complex
action to be taken by the mechanism, so that the pair of buttons could well be described
as "triggering" the song which is played. In the process which converts genotype into
phenotype, cellular jukeboxes-if you will pardon the notion!-accept "button-pushings"
from short excerpts from a long strand of DNA, and the "songs" which they play are
often prime ingredients in the creation of further "jukeboxes". It is as if the output of real
jukeboxes, instead of being love ballads, were songs whose lyrics told how to build more
complex jukeboxes ... Portions of the DNA trigger the manufacture of proteins; those
proteins trigger hundreds of new reactions; they in turn trigger the replicating-operation
which, in several steps, copies the DNA-and on and on ... This gives a sense of how
recursive the whole process is. The final result of these many-triggered triggerings is the
phenotype-the individual. And one says that the phenotype is the revelation-the "pulling-
out"-of the information that was present in the DNA to start with, latently. (The term
"revelation" in this context is due to




The Location of Meaning                                                                   168

Jacques Monod, one of the deepest and most original of twentieth-century molecular
biologists.)
        Now no one would say that a song coming out of the loudspeaker of jukebox
constitutes a "revelation" of information inherent in the pair buttons which were pressed,
for the pair of buttons seem to be mere triggers, whose purpose is to activate information-
bearing portions of the jukebox mechanism. On the other hand, it seems perfectly
reasonable to call t extraction of music from a record a "revelation" of information
inherent the record, for several reasons:

  (1) the music does not seem to be concealed in the mechanism of the record player;
  (2) it is possible to match pieces of the input (the record) with pieces of the output
      (the music) to an arbitrary degree of accuracy;
  (3) it is possible to play other records on the same record player and get other
      sounds out;
  (4) the record and the record player are easily separated from one another.

It is another question altogether whether the fragments of a smashed record contain
intrinsic meaning. The edges of the separate pieces together and in that way allow the
information to be reconstituted-t something much more complex is going on here. Then
there is the question of the intrinsic meaning of a scrambled telephone call ... There is a
vast spectrum of degrees of inherency of meaning. It is interesting to try place epigenesis
in this spectrum. As development of an organism takes place, can it be said that the
information is being "pulled out" of its DNA? Is that where all of the information about
the organism's structure reside;

                  DNA and the Necessity of Chemical Context

In one sense, the answer seems to be yes, thanks to experiments li Avery's. But in another
sense, the answer seems to be no, because so much of the pulling-out process depends on
extraordinarily complicated cellular chemical processes, which are not coded for in the
DNA itself. The DNA relies on the fact that they will happen, but does not seem to
contain a code which brings them about. Thus we have two conflicting views on the
nature of the information in a genotype. One view says that so much of t information is
outside the DNA that it is not reasonable to look upon the DNA as anything more than a
very intricate set of triggers, like a sequence of buttons to be pushed on a jukebox;
another view says that the information is all there, but in a very implicit form.
        Now it might seem that these are just two ways of saying the same thing, but that
is not necessarily so. One view says that the DNA is quite meaningless out of context; the
other says that even if it were taken out context, a molecule of DNA from a living being
has such a compelling inner




The Location of Meaning                                                                169

logic to its structure that its message could be deduced anyway. To put it as succinctly as
possible, one view says that in order for DNA to have meaning, chemical context is
necessary; the other view says that only intelligence is necessary to reveal the "intrinsic
meaning" of a strand of DNA.

                                    An Unlikely UFO
        We can get some perspective on this issue by considering a strange hypothetical
event. A record of David Oistrakh and Lev Oborin playing Bach's sonata in F Minor for
violin and clavier is sent up in a satellite. From the satellite it is then launched on a course
which will carry it outside of the solar system, perhaps out of the entire galaxy just a thin
plastic platter with a hole in the middle, swirling its way through intergalactic space. It
has certainly lost its context. How much meaning does it carry?
        If an alien civilization were to encounter it, they would almost certainly be struck
by its shape, and would probably be very interested in it. Thus immediately its shape,
acting as a trigger, has given them some information: that it is an artifact, perhaps an
information-bearing artifact. This idea-communicated, or triggered, by the record itself-
now creates a new context in which the record will henceforth be perceived. The next
steps in the decoding might take considerably longer-but that is very hard for us to assess.
We can imagine that if such a record had arrived on earth in Bach's time, no one would
have known what to make of it, and very likely it would not have gotten deciphered. But
that does not diminish our conviction that the information was in principle there; we just
know that human knowledge in those times was not very sophisticated with respect to the
possibilities of storage, transformation, and revelation of information.

                      Levels of Understanding of a Message
Nowadays, the idea of decoding is extremely widespread; it is a significant part of the
activity of astronomers, linguists, archaeologists, military specialists, and so on. It is
often suggested that we may be floating in a sea of radio messages from other
civilizations, messages which we do not yet know how to decipher. And much serious
thought has been given to the techniques of deciphering such a message. One of the main
problems perhaps the deepest problem-is the question, "How will we recognize the fact
that there is a message at all? How to identify a frame?" The sending of a record seems to
be a simple solution-its gross physical structure is very attention-drawing, and it is at
least plausible to us that it would trigger, in any sufficiently great intelligence, the idea of
looking for information hidden in it. However, for technological reasons, sending of solid
objects to other star systems seems to be out of the question. Still, that does not prevent
our thinking about the idea.
        Now suppose that an alien civilization hit upon the idea that the appropriate
mechanism for translation of the record is a machine which




The Location of Meaning                                                                    170

converts the groove-patterns into sounds. This would still be a far cry from a true
deciphering. What, indeed, would constitute a successful deciphering of such a record?
Evidently, the civilization would have to be able to ma sense out of the sounds. Mere
production of sounds is in itself hart worthwhile, unless they have the desired triggering
effect in the brains that is the word) of the alien creatures. And what is that desired
effect? would be to activate structures in their brains which create emotional effects in
them which are analogous to the emotional effects which experience in hearing the piece.
In fact, the production of sounds cot even be bypassed, provided that they used the record
in some other way get at the appropriate structures in their brains. (If we humans had a w
of triggering the appropriate structures in our brains in sequential order, as music does,
we might be quite content to bypass the sounds-but it see] extraordinarily unlikely that
there is any way to do that, other than via o ears. Deaf composers-Beethoven, Dvofák,
Faure-or musicians who can "hear" music by looking at a score, do not give the lie to this
assertion, for such abilities are founded upon preceding decades of direct auditory
experiences.)
         Here is where things become very unclear. Will beings of an alien civilization
have emotions? Will their emotions-supposing they have some-be mappable, in any
sense, onto ours? If they do have emotions somewhat like ours, do the emotions cluster
together in somewhat the same way as ours do? Will they understand such amalgams as
tragic beauty courageous suffering? If it turns out that beings throughout the universe do
share cognitive structures with us to the extent that even emotions overlap, then in some
sense, the record can never be out of its natural context; that context is part of the scheme
of things, in nature. And if such is the case, then it is likely that a meandering record, if
not destroyed en route, would eventually get picked up by a being or group of beings, at
get deciphered in a way which we would consider successful.


                               "Imaginary Spacescape"
        In asking about the meaning of a molecule of DNA above, I used t phrase
"compelling inner logic"; and I think this is a key notion. To illustrate this, let us slightly
modify our hypothetical record-into-spa event by substituting John Cage's "Imaginary
Landscape no. 4" for the Bach. This piece is a classic of aleatoric, or chance, music-
music who structure is chosen by various random processes, rather than by an attempt to
convey a personal emotion. In this case, twenty-four performers attar themselves to the
twenty-four knobs on twelve radios. For the duration the piece they twiddle their knobs in
aleatoric ways so that each radio randomly gets louder and softer, switching stations all
the while. The tot sound produced is the piece of music. Cage's attitude is expressed in 14
own words: "to let sounds be themselves, rather than vehicles for man made theories or
expressions of human sentiments."




The Location of Meaning                                                                    171

        Now imagine that this is the piece on the record sent out into space. It would be
extraordinarily unlikely-if not downright impossible-for an alien civilization to
understand the nature of the artifact. They would probably be very puzzled by the
contradiction between the frame message ("I am a message; decode me"), and the chaos
of the inner structure. There are few "chunks" to seize onto in this Cage piece, few
patterns which could guide a decipherer. On the other hand, there seems to be, in a Bach
piece, much to seize onto-patterns, patterns of patterns, and so on. We have no way of
knowing whether such patterns are universally appealing. We do not know enough about
the nature of intelligence, emotions, or music to say whether the inner logic of a piece by
Bach is so universally compelling that its meaning could span galaxies.
        However, whether Bach in particular has enough inner logic is not the issue here;
the issue is whether any message has, per se, enough compelling inner logic that its
context will be restored automatically whenever intelligence of a high enough level
comes in contact with it. If some message did have that context-restoring property, then it
would seem reasonable to consider the meaning of the message as an inherent property of
the message.

                              The Heroic Decipherers
Another illuminating example of these ideas is the decipherment of ancient texts written
in unknown languages and unknown alphabets. The intuition feels that there is
information inherent in such texts, whether or not we succeed in revealing it. It is as
strong a feeling as the belief that there is meaning inherent in a newspaper written in
Chinese, even if we are completely ignorant of Chinese. Once the script or language of a
text has been broken, then no one questions where the meaning resides: clearly it resides
in the text, not in the method of decipherment just as music resides in a record, not inside
a record player! One of the ways that we identify decoding mechanisms is by the fact that
they do not add any meaning to the signs or objects which they take as input; they merely
reveal the intrinsic meaning of those signs or objects. A jukebox is not a decoding
mechanism, for it does not reveal any meaning belonging to its input symbols; on the
contrary, it supplies meaning concealed inside itself.
         Now the decipherment of an ancient text may have involved decades of labor by
several rival teams of scholars, drawing on knowledge stored in libraries all over the
world ... Doesn't this process add information, too? Just how intrinsic is the meaning of a
text, when such mammoth efforts are required in order to find the decoding rules? Has
one put meaning into the text, or was that meaning already there? My intuition says that
the meaning was always there, and that despite the arduousness of the pulling-out
process, no meaning was pulled out that wasn't in the text to start with. This intuition
comes mainly from one fact: I feel that the result was inevitable; that, had the text not
been deciphered by this group at this time, it would have been deciphered by that group
at that time-and it would have come




The Location of Meaning                                                                 172

               FIGURE 39. The Rosetta Stone [courtesy of the British Museum.

out the same way. That is why the meaning is part of the text itself; it acts upon
intelligence in a predictable way. Generally, we can say: meaning is part of an object to
the extent that it acts upon intelligence in a predictable way.
        In Figure 39 is shown the Rosetta stone, one of the most precious of all historic
discoveries. It was the key to the decipherment of Egyptian hieroglyphics, for it contains
parallel text in three ancient scripts: hieroglyphic demotic characters, and Greek. The
inscription on this basalt stele was firs deciphered in 1821 by Jean Francois Champollion,
the "father of Egyptology"; it is a decree of priests assembled at Memphis in favor of
Ptolemy Epiphanes.




The Location of Meaning                                                               173

Three Layers of Any Message
        In these examples of decipherment of out-of-context messages, we can separate
out fairly clearly three levels of information: (1) the frame message; (2) the outer
message; (3) the inner message. The one we are most familiar with is (3), the inner
message; it is the message which is supposed to be transmitted: the emotional
experiences in music, the phenotype in genetics, the royalty and rites of ancient
civilizations in tablets, etc.

    To understand the inner message is to have extracted the meaning intended by the
                                        sender..

        The frame message is the message "I am a message; decode me if you can!"; and
it is implicitly conveyed by the gross structural aspects of any information-bearer.

        To understand the frame message is to recognize the need for a decoding-
                                     mechanism.

        If the frame message is recognized as such, then attention is switched to level (2),
the outer message. This is information, implicitly carried by symbol-patterns and
structures in the message, which tells how to decode the inner message.

     To understand the outer message is to build, or know how to build, the correct
                     decoding mechanism for the inner message.

This outer level is perforce an implicit message, in the sense that the sender cannot ensure
that it will be understood. It would be a vain effort to send instructions which tell how to
decode the outer message, for they would have to be part of the inner message, which can
only be understood once the decoding mechanism has been found. For this reason, the
outer message is necessarily a set of triggers, rather than a message which can be
revealed by a known decoder.
         The formulation of these three "layers" is only a rather crude beginning at
analyzing how meaning is contained in messages. There may be layers and layers of
outer and inner messages, rather than just one of each. Think, for instance, of how
intricately tangled are the inner and outer messages of the Rosetta stone. To decode a
message fully, one would have to reconstruct the entire semantic structure which
underlay its creation and thus to understand the sender in every deep way. Hence one
could throw away the inner message, because if one truly understood all the finesses of
the outer message, the inner message would be reconstructible.
         The book After Babel, by George Steiner, is a long discussion of the interaction
between inner and outer messages (though he never uses that terminology). The tone of
his book is given by this quote:

   We normally use a shorthand beneath which there lies a wealth of subconscious,
   deliberately concealed or declared associations so extensive and intri-




The Location of Meaning                                                                 174

   cate that they probably equal the sum and uniqueness of our status as an individual person.'

Thoughts along the same lines are expressed by Leonard B. Meyer, in h book Music, the
Arts, and Ideas:

   The way of listening to a composition by Elliott Carter is radically different from the way
   of listening appropriate to a work by John Cage. Similarly, a novel by Beckett must in a
   significant sense be read differently from one by Bellow. A painting by Willem de
   Kooning and one by Andy Warhol require different perceptional-cognitive attitudes.'

        Perhaps works of art are trying to convey their style more than an thing else. In
that case, if you could ever plumb a style to its very bottom you could dispense with the
creations in that style. "Style", "outer message "decoding technique"-all ways of
expressing the same basic idea.

                          Schrodinger's Aperiodic Crystals
        What makes us see a frame message in certain objects, but none in other; Why
should an alien civilization suspect, if they intercept an errant record that a message lurks
within? What would make a record any different from a meteorite? Clearly its geometric
shape is the first clue that "something funny is going on". The next clue is that, on a more
microscopic scale, consists of a very long aperiodic sequence of patterns, arranged in a
spiral If we were to unwrap the spiral, we would have one huge linear sequence (around
2000 feet long) of minuscule symbols. This is not so different from a DNA molecule,
whose symbols, drawn from a meager "alphabet" of four different chemical bases, are
arrayed in a one-dimensional sequence, an then coiled up into a helix. Before Avery had
established the connection between genes and DNA, the physicist Erwin Schrödinger
predicted, o purely theoretical grounds, that genetic information would have to be stored
in "aperiodic crystals", in his influential book What Is Life? In fact books themselves are
aperiodic crystals contained inside neat geometric forms. These examples suggest that,
where an aperiodic crystal is found "packaged" inside a very regular geometric structure,
there may lurk a inner message. (I don't claim this is a complete characterization of frame
messages; however, it is a fact that many common messages have frame messages of this
description. See Figure 40 for some good examples.)

                           Languages for the Three Levels
The three levels are very clear in the case of a message found in a bottle washed up on a
beach. The first level, the frame message, is found when one picks up the bottle and sees
that it is sealed, and contains a dry piece c paper. Even without seeing writing, one
recognizes this type of artifact an information-bearer, and at this point it would take an
extraordinary almost inhuman-lack of curiosity, to drop the bottle and not look further.




The Location of Meaning                                                                       175

                          '




The Location of Meaning       176

Next, one opens the bottle and examines the marks on the paper. Perhaps, they are in
Japanese; this can be discovered without any of the inner message being understood-it
merely comes from a recognition of 1 characters. The outer message can be stated as an
English sentence: "I in Japanese." Once this has been discovered, then one can proceed
the inner message, which may be a call for help, a haiku poem, a lover’s lament ...
         It would be of no use to include in the inner message a translation the sentence
"This message is in Japanese", since it would take someone who knew Japanese to read
it. And before reading it, he would have recognize the fact that, as it is in Japanese, he
can read it. You might try wriggle out of this by including translations of the statement
"This mess2 is in Japanese" into many different languages. That would help it practical
sense, but in a theoretical sense the same difficulty is there. . English-speaking person
still has to recognize the "Englishness" of the message; otherwise it does no good. Thus
one cannot avoid the problem that one has to find out how to decipher the inner message
from the outside the inner message itself may provide clues and confirmations, but those ;
at best triggers acting upon the bottle finder (or upon the people whom enlists to help).
         Similar kinds of problem confront the shortwave radio listener. First he has to
decide whether the sounds he hears actually constitute a message or are just static. The
sounds in themselves do not give the answer, not e\% in the unlikely case that the inner
message is in the listener's own native language, and is saying, "These sounds actually
constitute a message a are not just static!" If the listener recognizes a frame message in
the soup then he tries to identify the language the broadcast is in-and clearly, he is still on
the outside; he accepts triggers from the radio, but they cam explicitly tell him the
answer.
         It is in the nature of outer messages that they are not conveyed in any

FIGURE 40. A collage of scripts. Uppermost on the left is an inscription in the un ciphered
boustrophedonic writing system from Easter Island, in which every second lin upside down. The
characters are chiseled on a wooden tablet, 4 inches by 35 inches. Mov clockwise, we encounter
vertically written Mongolian: above, present-day Mongolian, below, a document dating from
1314. Then we come to a poem in Bengali by Rabindran Tagore in the bottom righthand corner.
Next to it is a newspaper headline in Malayalam (II Kerala, southern India), above which is the
elegant curvilinear language Tamil (F Kerala). The smallest entry is part of a folk tale in
Buginese (Celebes Island, Indonesia). In center of the collage is a paragraph in the Thai
language, and above it a manuscript in Rn dating from the fourteenth century, containing a
sample of the provincial law of Scania (so Sweden). Finally, wedged in on the left is a section of
the laws of Hammurabi, written Assyrian cuneiform. As an outsider, I feel a deep sense of
mystery as I wonder how meanin cloaked in the strange curves and angles of each of these
beautiful aperiodic crystals. Info there is content. [From Ham Jensen, Sign, Symbol, and Script
(New York: G. Putnam's S. 1969), pp. 89 (cuneiform), 356 (Easter Island), 386, 417 (Mongolian),
552 (Runic); from Keno Katzner, The Languages of the World (New York: Funk \& Wagnalls,
1975), pp. 190 (Bengali),
         (Buginese); from I. A. Richards and Christine Gibson, English Through Pictures (New Y
Washington Square Press, 1960), pp. 73 (Tamil), 82 (Thai).




The Location of Meaning                                                                      177

explicit language. To find an explicit language in which to convey outer messages would
not be a breakthrough-it would be a contradiction in terms! It is always the listener's
burden to understand the outer message. Success lets him break through into the inside, at
which point the ratio of triggers to explicit meanings shifts drastically towards the latter.
By comparison with the previous stages, understanding the inner message seems
effortless. It is as if it just gets pumped in.

                       The "Jukebox" Theory of Meaning.
These examples may appear to be evidence for the viewpoint that no message has
intrinsic meaning, for in order to understand any inner message, no matter how simple it
is, one must first understand its frame message and its outer message, both of which are
carried only by triggers (such as being written in the Japanese alphabet, or having
spiraling grooves, etc.). It begins to seem, then, that one cannot get away from a
"jukebox" theory of meaning-the doctrine that no message contains inherent meaning,
because, before any message can be understood, it has to be used as the input to some
"jukebox", which means that information contained in the "jukebox" must be added to the
message before it acquires meaning.
        This argument is very similar to the trap which the Tortoise caught Achilles in, in
Lewis Carroll's Dialogue. There, the trap was the idea that before you can use any rule,
you have to have a rule which tells you how to use that rule; in other words, there is an
infinite hierarchy of levels of rules, which prevents any rule from ever getting used. Here,
the trap is the idea that before you can understand any message, you have to have a
message which tells you how to understand that message; in other words, there is an
infinite hierarchy of levels of messages, which prevents any message from ever getting
understood. However, we all know that these paradoxes are invalid, for rules do get used,
and messages do get understood. How come?

Against the Jukebox Theory
This happens because our intelligence is not disembodied, but is instantiated in physical
objects: our brains. Their structure is due to the long process of evolution, and their
operations are governed by the laws of physics. Since they are physical entities, our
brains run without being told how to run. So it is at the level where thoughts are produced
by physical law that Carroll's rule-paradox breaks down; and likewise, it is at the level
where a brain interprets incoming data as a message that the message-paradox breaks
down. It seems that brains come equipped with "hardware" for recognizing that certain
things are messages, and for decoding those messages. This minimal inborn ability to
extract inner meaning is what allows the highly recursive, snowballing process of
language acquisition to take place. The inborn hardware is like a jukebox: it supplies the
additional information which turns mere triggers into complete messages.




The Location of Meaning                                                                  178

                 Meaning Is Intrinsic If Intelligence Is Natural
Now if different people's "jukeboxes" had different "songs" in then responded to given
triggers in completely idiosyncratic ways, the would have no inclination to attribute
intrinsic meaning to those tri; However, human brains are so constructed that one brain
responds in much the same way to a given trigger as does another brain, all other t being
equal. This is why a baby can learn any language; it responds to triggers in the same way
as any other baby. This uniformity of "human jukeboxes" establishes a uniform
"language" in which frame message outer messages can be communicated. If,
furthermore, we believe human intelligence is just one example of a general phenomena
nature-the emergence of intelligent beings in widely varying contexts then presumably
the "language" in which frame messages and outer sages are communicated among
humans is a "dialect" of a universal gauge by which intelligences can communicate with
each other. Thus, would be certain kinds of triggers which would have "universal
triggering power", in that all intelligent beings would tend to respond to them i same way
as we do.
        This would allow us to shift our description of where meaning located. We could
ascribe the meanings (frame, outer, and inner) message to the message itself, because of
the fact that deciphering mechanisms are themselves universal-that is, they are
fundamental f of nature which arise in the same way in diverse contexts. To make it
concrete, suppose that "A-5" triggered the same song in all jukeboxes suppose moreover
that jukeboxes were not man-made artifacts, but w occurring natural objects, like galaxies
or carbon atoms. Under such circumstances, we would probably feel justified in calling
the universal triggering power of "A-5" its "inherent meaning"; also, "A-5" would merit:
the name of "message", rather than "trigger", and the song would indeed "revelation" of
the inherent, though implicit, meaning of "A-5".

                                    Earth Chauvinism
This ascribing of meaning to a message comes from the invariance c processing of the
message by intelligences distributed anywhere ii universe. In that sense, it bears some
resemblance to the ascribing of to an object. To the ancients, it must have seemed that an
object's weight was an intrinsic property of the object. But as gravity became understood,
it was realized that weight varies with the gravitational field the object is immersed in.
Nevertheless, there is a related quantity, the mass, which not vary according to the
gravitational field; and from this invariance the conclusion that an object's mass was an
intrinsic property of the object itself. If it turns out that mass is also variable, according to
context, then will backtrack and revise our opinion that it is an intrinsic property of an
object. In the same way, we might imagine that there could exist other




The Location of Meaning                                                                     179

kinds of "jukeboxes"-intelligences-which communicate among each other via messages
which we would never recognize as messages, and who also would never recognize our
messages as messages. If that were the case, then the claim that meaning is an intrinsic
property of a set of symbols would have to be reconsidered. On the other hand, how
could we ever realize that such beings existed?
        It is interesting to compare this argument for the inherency of meaning with a
parallel argument for the inherency of weight. Suppose one defined an object's weight as
"the magnitude of the downward force which the object exerts when on the surface of the
planet Earth". Under this definition, the downward force which an object exerts when on
the surface of Mars would have to be given another name than "weight". This definition
makes weight an inherent property, but at the cost of geocentricity" Earth chauvinism". It
would be like "Greenwich chauvinism"-refusing to accept local time anywhere on the
globe but in the GMT time zone. It is an unnatural way to think of time.
        Perhaps we are unknowingly burdened with a similar chauvinism with respect to
intelligence, and consequently with respect to meaning. In our chauvinism, we would call
any being with a brain sufficiently much like our own "intelligent", and refuse to
recognize other types of objects as intelligent. To take an extreme example, consider a
meteorite which, instead of deciphering the outer-space Bach record, punctures it with
colossal indifference, and continues in its merry orbit. It has interacted with the record in
a way which we feel disregards the record's meaning. Therefore, we might well feel
tempted to call the meteorite "stupid". But perhaps we would thereby do the meteorite a
disservice. Perhaps it has a "higher intelligence" which we in our Earth chauvinism
cannot perceive, and its interaction with the record was a manifestation of that higher
intelligence. Perhaps, then, the record has a "higher meaning"-totally different from that
which we attribute to it; perhaps its meaning depends on the type of intelligence
perceiving it. Perhaps.
        It would be nice if we could define intelligence in some other way than "that
which gets the same meaning out of a sequence of symbols as we do". For if we can only
define it this one way, then our argument that meaning is an intrinsic property is circular,
hence content-free. We should try to formulate in some independent way a set of
characteristics which deserve the name "intelligence". Such characteristics would
constitute the uniform core of intelligence, shared by humans. At this point in history we
do not yet have a well-defined list of those characteristics. However, it appears likely that
within the next few decades there will be much progress made in elucidating what human
intelligence is. In particular, perhaps cognitive psychologists, workers in Artificial
Intelligence, and neuroscientists will be able to synthesize their understandings, and come
up with a definition of intelligence. It may still be human-chauvinistic; there is no way
around that. But to counterbalance that, there may be some elegant and beautiful-and
perhaps even simple-abstract ways of characterizing the essence of intelligence. This
would serve to lessen the feeling of having




The Location of Meaning                                                                  180

formulated an anthropocentric concept. And of course, if contact were established with an
alien civilization from another star system, we feel supported in our belief that our own
type of intelligence is not just a fluke, but an example of a basic form which reappears in
nature in contexts, like stars and uranium nuclei. This in turn would support the idea of
meaning being an inherent property.
        To conclude this topic, let us consider some new and old ex; and discuss the
degree of inherent meaning which they have, by ourselves, to the extent that we can, in
the shoes of an alien civilization which intercepts a weird object ...


                                Two Plaques in Space
Consider a rectangular plaque made of an indestructible metallic alloy which are
engraved two dots, one immediately above the another preceding colon shows a picture.
Though the overall form of the might suggest that it is an artifact, and therefore that it
might conceal some message, two dots are simply not sufficient to convey anything. (Can
before reading on, hypothesize what they are supposed to mean suppose that we made a
second plaque, containing more dots, as follows.

                                  .
                                 ..
                                …
                                …..
                               ……..
                              ……………
                            ……………………..
                          …………………………………


         Now one of the most obvious things to do-so it might seer terrestrial intelligence
at least-would be to count the dots in the successive rows. The sequence obtained is:

                            1, 1, 2, 3, 5, 8, 13, 21, 34.


Here there is evidence of a rule governing the progression from one the next. In fact, the
recursive part of the definition of the Fib numbers can be inferred, with some confidence,
from this list. Supp think of the initial pair of values (1,1) as a "genotype" from which the
"phenotype"-the full Fibonacci sequence-is pulled out by a recursive rule. By sending the
genotype alone-namely the first version plaque-we fail to send the information which
allows reconstitution phenotype. Thus, the genotype does not contain the full
specification of



The Location of Meaning                                                                  181

the phenotype. On the other hand, if we consider the second version of the plaque to be
the genotype, then there is much better cause to suppose that the phenotype could
actually be reconstituted. This new version of the genotype-a "long genotype"-contains so
much information that the mechanism by which phenotype is pulled out of genotype can
be inferred by intelligence from the genotype alone.
        Once this mechanism is firmly established as the way to pull phenotype from
genotype, then we can go back to using "short genotypes"-like the first plaque. For
instance, the "short genotype" (1,3) would yield the phenotype

                1,      3,     4,      7,      11,     18,     29,     47, .. .

-the Lucas sequence. And for every set of two initial values-that is, for every short
genotype-there will be a corresponding phenotype. But the short genotypes, unlike the
long ones, are only triggers-buttons to be pushed on the jukeboxes into which the
recursive rule has been built. The long genotypes are informative enough that they
trigger, in an intelligent being, the recognition of what kind of "jukebox" to build. In that
sense, the long genotypes contain the information of the phenotype, whereas the short
genotypes do not. In other words, the long genotype transmits not only an inner message,
but also an outer message, which enables the inner message to be read. It seems that the
clarity of the outer message resides in the sheer length of the message. This is not
unexpected; it parallels precisely what happens in deciphering ancient texts. Clearly,
one's likelihood of success depends crucially on the amount of text available.

                                 Bach vs. Cage Again
But just having a long text may not be enough. Let us take up once more the difference
between sending a record of Bach's music into space, and a record of John Cage's music.
Incidentally, the latter, being a Composition of Aleatorically Generated Elements, might
be handily called a "CAGE", whereas the former, being a Beautiful Aperiodic Crystal of
Harmony, might aptly be dubbed a "BACH". Now let's consider what the meaning of a
Cage piece is to ourselves. A Cage piece has to be taken in a large cultural setting-as a
revolt against certain kinds of traditions. Thus, if we want to transmit that meaning, we
must not only send the notes of the piece, but we must have earlier communicated an
extensive history of Western culture. It is fair to say, then, that an isolated record of John
Cage's music does not have an intrinsic meaning. However, for a listener who is
sufficiently well versed in Western and Eastern cultures, particularly in the trends in
Western music over the last few decades, it does carry meaning-but such a listener is like
a jukebox, and the piece is like a pair of buttons. The meaning is mostly contained inside
the listener to begin with; the music serves only to trigger it. And this "jukebox", unlike
pure intelligence, is not at all universal; it is highly earthbound, depending on
idiosyncratic se-




The Location of Meaning                                                                   182

quences of events all over our globe for long period of time. Hoping that John Cage's
music will be understood by another civilization is like hoping that your favorite tune, on
a jukebox on the moon, will have the same buttons as in a saloon in Saskatoon.
        On the other hand, to appreciate Bach requires far less cultural k edge. This may
seem like high irony, for Bach is so much more con and organized, and Cage is so devoid
of intellectuality. But there strange reversal here: intelligence loves patterns and balks at
randomness For most people, the randomness in Cage's music requires much explanation;
and even after explanations, they may feel they are missing the message-whereas with
much of Bach, words are superfluous. In sense, Bach's music is more self-contained than
Cage's music. Still, it is clear how much of the human condition is presumed by Bach.
        For instance, music has three major dimensions of structure (me harmony,
rhythm), each of which can be further divided into small intermediate, and overall
aspects. Now in each of these dimensions, there is a certain amount of complexity which
our minds can handle before boggling; clearly a composer takes this into account, mostly
unconsciously when writing a piece. These "levels of tolerable complexity" along
different dimensions are probably very dependent on the peculiar conditions of our
evolution as a species, and another intelligent species might have developed music with
totally different levels of tolerable complexity along these many dimensions. Thus a Bach
piece might conceivably have to be accompanied, by a lot of information about the
human species, which simply could not inferred from the music's structure alone. If we
equate the Bach music a genotype, and the emotions which it is supposed to evoke with
the phenotype, then what we are interested in is whether the genotype con all the
information necessary for the revelation of the phenotype.

How Universal Is DNA's Message?
The general question which we are facing, and which is very similar t questions inspired
by the two plaques, is this: "How much of the co necessary for its own understanding is a
message capable of restoring? can now revert to the original biological meanings of
"genotype" "phenotype"-DNA and a living organism-and ask similar quest Does DNA
have universal triggering power? Or does it need a "biojukebox" to reveal its meaning?
Can DNA evoke a phenotype without being embedded in the proper chemical context?
To this question to answer is no-but a qualified no. Certainly a molecule of DNA in a
vacuum will not create anything at all. However, if a molecule of DNA were set to seek
its fortune in the universe, as we imagined the BACH and the CAGE were, it might be
intercepted by an intelligent civilization. They might first of all recognize its frame
message. Given that, they might to try to deduce from its chemical structure what kind of
chemical environment it seemed to want, and then supply such an environment. Succes-




The Location of Meaning                                                                  183

sively more refined attempts along these lines might eventually lead to a full restoration
of the chemical context necessary for the revelation of DNA's phenotypical meaning.
This may sound a little implausible, but if one allows many millions of years for the
experiment, perhaps the DNA's meaning would finally emerge.
        On the other hand, if the sequence of bases which compose a strand of DNA were
sent as abstract symbols (as in Fig. 41), not as a long helical molecule, the odds are
virtually nil that this, as an outer message, would trigger the proper decoding mechanism
which would enable the phenotype to be drawn out of the genotype. This would be a case
of wrapping an inner message in such an abstract outer message that the context-restoring
power of the outer message would be lost, and so in a very pragmatic sense, the set of
symbols would have no intrinsic meaning. Lest you think this all sounds hopelessly
abstract and philosophical, consider that the exact moment when phenotype can be said to
be "available", or "implied", by genotype, is a highly charged issue in our day: it is the
issue of abortion.

FIGURE 41. This Giant Aperiodic Crystal is the base sequence for the chromosome of
bacteriophage OX174. It is the first complete genome ever mapped out for any organism.
About 2,000 of these boustrophedonic pages would be needed to show the base sequence
of a single E. Coli cell, and about one million pages to show the base sequence of the
DNA of a single human cell. The book now in your hands contains roughly the same
amount of information as a molecular blueprint for one measly E. Coli cell.




The Location of Meaning                                                               184

                  Chromatic Fantasy, And Feud.
       Having had a splendid dip in the pond, the Tortoise is just crawling out and
                 shaking himself dry, when who but Achilles walks by.

Tortoise: Ho there, Achilles. I was just thinking of you as I splash around in the pond.
Achilles: Isn't that curious? I was just thinking of you, too, while I meandered through the
       meadows. They're so green at this time of year.
Tortoise: You think so? It reminds me of a thought I was hoping to share with you.
       Would you like to hear it?
Achilles: Oh, I would be delighted. That is, I would be delighted as long you're not going
       to try to snare me in one of your wicked traps of log Mr. T.
Tortoise: Wicked traps? Oh, you do me wrong. Would I do anything wicked? I'm a
       peaceful soul, bothering nobody and leading a gent; herbivorous life. And my
       thoughts merely drift among the oddities and quirks of how things are (as I see
       them). I, humble observer phenomena, plod along and puff my silly words into
       the air rather unspectacularly, I am afraid. But to reassure you about my intention
       I was only planning to speak of my Tortoise-shell today, and as you know, those
       things have nothing-nothing whatsoever-to do with logic!
Achilles: Your words Do reassure me, Mr. T. And, in fact, my curiosity quite piqued. I
       would certainly like to listen to what you have to say even if it is unspectacular.
Tortoise: Let's see ... how shall I begin? Hmm ... What strikes you me about my shell,
       Achilles?
Achilles: It looks wonderfully clean!
Tortoise: Thank you. I just went swimming and washed off several layers of dirt which
       had accumulated last century. Now you can see ho green my shell is.
Achilles: Such a good healthy green shell, it's nice to see it shining in sun.
Tortoise: Green? It's not green.
Achilles: Well, didn't you just tell me Tortoise: I did.
Achilles: Then, we agree: it is green. Tortoise: No, it isn't green.
Achilles: Oh, I understand your game. You're hinting to me that what you say isn't
       necessarily true; that Tortoises play with language; that your statements and
       reality don't necessarily match; that --

Tortoise: I certainly am not. Tortoises treat words as sacred. Tortoises revere accuracy.
Achilles: Well, then, why did you say that your shell is green, and that it is not green
       also?
Tortoise: I never said such a thing; but I wish I had. Achilles: You would have liked to
       say that?
Tortoise: Not a bit. I regret saying it, and disagree wholeheartedly with it. Achilles: That
       certainly contradicts what you said before!
Tortoise: Contradicts? Contradicts? I never contradict myself. It's not part of Tortoise-
       nature.
Achilles: Well, I've caught you this time, you slippery fellow, you. Caught you in a full-
       fledged contradiction.
Tortoise: Yes, I guess you did.
Achilles: There you go again! Now you're contradicting yourself more and more! You are
       so steeped in contradiction it's impossible to argue with you!
Tortoise: Not really. I argue with myself without any trouble at all. Perhaps the problem
       is with you. I would venture a guess that maybe you're the one who's
       contradictory, but you're so trapped in your own tangled web that you can't see
       how inconsistent you're being.
Achilles: What an insulting suggestion! I'm going to show you that you're the
       contradictory one, and there are no two ways about it.
Tortoise: Well, if it's so, your task ought to be cut out for you. What could be easier than
       to point out a contradiction? Go ahead-try it out.
Achilles: Hmm ... Now I hardly know where to begin. Oh ... I know. You first said that
       (1) your shell is green, and then you went on to say that (2) your shell is not
       green. What more can I say?
Tortoise: Just kindly point out the contradiction. Quit beating around the bush.
Achilles: But-but-but ... Oh, now I begin to see. (Sometimes I am so slow-witted!) It must
       be that you and I differ as to what constitutes a contradiction. That's the trouble.
       Well, let me make myself very clear: a contradiction occurs when somebody says
       one thing and denies it at the same time.
Tortoise: A neat trick. I'd like to see it done. Probably ventriloquists would excel at
       contradictions, speaking out of both sides of their mouth, as it were. But I'm not a
       ventriloquist.
Achilles: Well, what I actually meant is just that somebody can say one thing and deny it
       all within one single sentence! It doesn't literally have to be in the same instant.
Tortoise: Well, you didn't give ONE sentence. You gave TWO.
Achilles: Yes-two sentences that contradict each other!
Tortoise: I am sad to see the tangled structure of your thoughts becoming so exposed,
       Achilles. First you told me that a contradiction is some thing which occurs in a
       single sentence. Then you told me that you

Found a contradiction in a pair of sentences I uttered. Frankly, it’s just as I said. Your
       own system of thought is so delusional that you manage to avoid seeing how
       inconsistent it is. From the outside, however plain as day.
Achilles: Sometimes I get so confused by your diversionary tactics tl can't quite tell if
       we're arguing about something utterly petty, or something deep and profound!
Tortoise: I assure you, Tortoises don't spend their time on the petty. Hence it's the latter.
Achilles: I am very reassured. Thank you. Now I have had a moment to reflect, and I see
       the necessary logical step to convince you that you contradicted yourself.
Tortoise: Good, good. I hope it's an easy step, an indisputable one.
Achilles: It certainly is. Even you will agree with it. The idea is that you believed
       sentence 1 ("My shell is green"), AND you believed sentence 2 ("My shell is not
       green"), you would believe one compound( sentence in which both were
       combined, wouldn't you?
Tortoise: Of course. It would only be reasonable ... providing just that the manner of
       combination is universally acceptable. But I'm sure we'll agree on that.
Achilles: Yes, and then I'll have you! The combination I propose is
Tortoise: But we must be careful in combining sentences. For instance you'd grant that
       "Politicians lie" is true, wouldn't you?
Achilles: Who could deny it?
Tortoise: Good. Likewise, "Cast-iron sinks" is a valid utterance, isn't it?
Achilles: Indubitably.
Tortoise: Then, putting them together, we get "Politicians lie in cast iron sinks". Now
       that's not the case, is it?
Achilles: Now wait a minute ... "Politicians lie in cast-iron sinks?" N no, but
Tortoise: So, you see, combining two true sentences in one is not a policy, is it?
Achilles: But you-you combined the two-in such a silly way!
Tortoise: Silly? What have you got to object to in the way I combined them Would you
       have me do otherwise?
Achilles: You should have used the word "and", not "in".
Tortoise: I should have? You mean, if YOU'D had YOUR way, I should h;
Achilles: No-it's the LOGICAL thing to do. It's got nothing to do with personally.
Tortoise: This is where you always lose me, when you resort to your L and its high-
       sounding Principles. None of that for me today, plea
Achilles: Oh, Mr. Tortoise, don't put me through all this agony. You k very well that
       that's what "and" means! It's harmless to combine true sentences with "and"!
Tortoise: "Harmless", my eye! What gall! This is certainly a pernicious plot

       to entrap a poor, innocent, bumbling Tortoise in a fatal contradiction. If it were so
       harmless, why would you be trying so bloody hard to get me to do it? Eh?
Achilles: You've left me speechless. You make me feel like a villain, where I really had
       only the most innocent of motivations.
Tortoise: That's what everyone believes of himself...
Achilles: Shame on me-trying to outwit you, to use words to snare you in a self-
       contradiction. I feel so rotten.
Tortoise: And well you should. I know what you were trying to set up. Your plan was to
       make me accept sentence 3, to wit: "My shell is green and my shell is not green".
       And such a blatant falsehood is repellent to the Tongue of a Tortoise.
Achilles: Oh, I'm so sorry I started all this.
Tortoise: You needn't be sorry. My feelings aren't hurt. After all, I'm used to the
       unreasonable ways of the folk about me. I enjoy your company, Achilles, even if
       your thinking lacks clarity.
Achilles: Yes ... Well, I fear I am set in my ways, and will probably continue to err and
       err again, in my quest for Truth.
Tortoise: Today's exchange may have served a little to right your course. Good day,
       Achilles.
Achilles: Good day, Mr. T.


                      The Propositional Calculus
                                  Words and Symbols
THE PRECEDING DIALOGUE is reminiscent of the Two-Part Invention by Lewis
Carroll. In both, the Tortoise refuses to use normal, ordinary in the normal, ordinary way-
or at least he refuses to do so when it is his advantage to do so. A way to think about the
Carroll paradox was given last Chapter. In this Chapter we are going to make symbols dc
Achilles couldn't make the Tortoise do with his words. That is, we are to make a formal
system one of whose symbols will do just what A wished the word `and' would do, when
spoken by the Tortoise, and ail of whose symbols will behave the way the words 'if... then
. . .' ought to behave. There are only two other words which we will attempt to deal with
`or' and `not'. Reasoning which depends only on correct usage of these words is termed
propositional reasoning.

            Alphabet and First Rule of the Propositional Calculus
I will present this new formal system, called the Propositional Calculus, like a puzzle, not
explaining everything at once, but letting you things out to some extent. We begin with
the list of symbols:

                                   <    >
                              P    Q    R      ´
                              ∧    ∨    ⊃     ~

The first rule of this system that I will reveal is the following:

RULE OF JOINING: If x and y are theorems of the system, then so is the string < x∧y >.

This rule takes two theorems and combines them into one. It s remind you of the
Dialogue.

Well-Formed Strings
There will be several other rules of inference, and they will all be pres shortly-but first, it
is important to define a subset of all strings, namely the




The Propositional Calculus                                                                 189

well formed strings. They will be defined in a recursive way. We begin with the
ATOMS: P, Q, and R are called atoms.. New atoms are formed by appending primes
onto the right of old atoms-thus, R', Q", P"', etc. This gives an endless supply of atoms.
All atoms are well-formed.

Then we have four recursive

FORMATION RULES: If x and y are well-formed, then the following four strings are
also well-formed:

(1)    ~x
(2)    < x∧y>
(3)    < x∨y>
(4)    < x⊃y>

For example, all of the following are well-formed:

P                                     atom
~P                                    by (1)
~~P                                   by (1)
Q´                                    atom
~Q1                                   by (1)
<P∧~Q' >                              by (2)
~<P∧~Q' >                             by (1)
~~<P⊃~Q' >                            by (4)
<~<P∧~Q' >∨~~<P⊃~Q' >>                by (3)

The last one may look quite formidable, but it is built up straightforwardly from two
components-namely the two lines just above it. Each of them is in turn built up from
previous lines ... and so on. Every well-formed string can in this way be traced back to its
elementary constituents-that is, atoms. You simply run the formation rules backwards
until you can no more. This process is guaranteed to terminate, since each formation rule
(when run forwards) is a lengthening rule, so that running it backwards always drives you
towards atoms.
        This method of decomposing strings thus serves as a check on the well-
formedness of any string. It is a top-down decision procedure for wellformedness. You
can test your understanding of this decision procedure by checking which of the
following strings are well-formed:

       (1) <P>
       (2) (2) <~P>
       (3) <P∧Q∧R>
       (4) <P∧Q>
       (5) <<P∧Q>∧Q~∧P>>
       (6) <P∧~P>
       (7) <<P∨<Q⊃R>>∧<~P∨~R´>>
       (8) <P∧Q>∧< Q∧P:

The Propositional Calculus                                                              190

(Answer: Those whose numbers are Fibonacci numbers are not formed. The rest are well-
formed.)

                               More Rules of Inference
Now we come to the rest of the rules by which theorems of this system constructed. A
few rules of inference follow. In all of them, the symbols ´x´ and 'y' are always to be
understood as restricted to well formed strings

RULE OF SEPARATION: If < x∧y> is a theorem, then both x and theorems.

Incidentally, you should have a pretty good guess by now as to concept the symbol `A'
stands for. (Hint: it is the troublesome word the preceding Dialogue.) From the following
rule, you should be a figure out what concept the tilde ('~') represents:

DOUBLE-TILDE RULE: The string '~~' can be deleted from any theorem. It can also be
inserted into any theorem, provided that the rest string is itself well-formed.


                                   The Fantasy Rule
Now a special feature of this system is that it has no axioms-only rule you think back to
the previous formal systems we've seen, you may w( how there can be any theorems,
then. How does everything get started? The answer is that there is one rule which
manufactures theorems from out of thin air-it doesn't need an "old theorem" as input.
(The rest of the do require input.) This special rule is called the fantasy rule. The reason I
call it that is quite simple.
         To use the fantasy rule, the first thing you do is to write down an well-formed
string x you like, and then "fantasize" by asking, "What if string x were an axiom, or a
theorem?" And then, you let the system give an answer. That is, you go ahead and make a
derivation with x ; opening line; let us suppose y is the last line. (Of course the derivation
must strictly follow the rules of the system.) Everything from x to y (inclusive) is the
fantasy; x is the premise of the fantasy, and y is its outcome. The next step is to jump out
of the fantasy, having learned from it that out.

                        If x were a theorem, y would be a theorem.

Still, you might wonder, where is the real theorem? The real theorem is the string

                                           <x⊃y>

Notice the resemblance of this string to the sentence printed above
       To signal the entry into, and emergence from, a fantasy, one uses the



The Propositional Calculus                                                                191

square brackets `[' and ']', respectively. Thus, whenever you see a left square bracket, you
know you are "pushing" into a fantasy, and the next line will contain the fantasy's
premise. Whenever you see a right square bracket, you know you are "popping" back out,
and the preceding line was the outcome. It is helpful (though not necessary) to indent
those lines of a derivation which take place in fantasies.
        Here is an illustration of the fantasy rule, in which the string P is taken as a
premise. (It so happens that P is not a theorem, but that is of no import; we are merely
inquiring, "What if it were?") We make the following fantasy:

       [                                     push into fantasy
           P                                 premise
           ~~~P                              outcome (by double tilde rule)
       ]                                     pop out of fantasy

       The fantasy shows that:

       If P were a theorem, so would ~~P be one.

       We now "squeeze" this sentence of English (the metalanguage) into the formal
notation (the object language): <P⊃~~P>. This, our first theorem of the Propositional
Calculus, should reveal to you the intended interpretation of the symbol `⊃'.
       Here is another derivation using the fantasy rule:

       [                                     push
           <P∧Q>                             premise
           P                                 separation
           Q                                 separation
           <Q∧P>                             joining
       ]                                     pop
       <<P∧Q>⊃<Q∧P>>                         fantasy rule

It is important to understand that only the last line is a genuine theorem, here-everything
else is in the fantasy.

                         Recursion and the Fantasy Rule
As you might guess from the recursion terminology "push" and "pop", the fantasy rule
can be used recursively-thus, there can be fantasies within fantasies, thrice-nested
fantasies, and so on. This means that there are all sorts of "levels of reality", just as in
nested stories or movies. When you pop out of a movie-within-a-movie, you feel for a
moment as if you had reached the real world, though you are still one level away from the
top. Similarly, when you pop out of a fantasy-within-a-fantasy, you are in a "realer"
world than you had been, but you are still one level away from the top.
        Now a "No Smoking" sign inside a movie theater does not apply to the



The Propositional Calculus                                                              192

characters in the movie-there is no carry-over from the real world in fantasy world, in
movies. But in the Propositional Calculus, then carry-over from the real world into the
fantasies; there is even carry from a fantasy to fantasies inside it. This is formalized by
the following rule:

CARRY-OVER RULE: Inside a fantasy, any theorem from the "reality level higher can
    be brought in and used.

It is as if a "No Smoking" sign in a theater applied not only to a moviegoers, but also to
all the actors in the movie, and, by repetition of the same idea, to anyone inside multiply
nested movies! (Warning: There carry-over in the reverse direction: theorems inside
fantasies cannot be exported to the exterior! If it weren't for this fact, you could write any
as the first line of a fantasy, and then lift it out into the real world as a theorem.)

        To show how carry-over works, and to show how the fantasy rule can be used
recursively, we present the following derivation:

       [                                      push
           P                                  premise of outer fantasy
           [                                  push again
               Q                              premise of inner fantasy
               P                              carry-over of P into inner fantasy
               <P∧Q>                          joining
           ]                                  pop out of inner fantasy, regain outer fantasy
           <Q⊃<P∧Q>>                          fantasy rule
       ]                                      pop out of outer fantasy, reach real world!
       <P⊃<Q⊃<P∧Q>>>                          fantasy rule

        Note that I've indented the outer fantasy once, and the inner fantasy twice, to
emphasize the nature of these nested "levels of reality". One to look at the fantasy rule is
to say that an observation made about the system is inserted into the system. Namely, the
theorem < x⊃y> which gets produced can be thought of as a representation inside the
system of the statement about the system "If x is a theorem, then y is too". To be specific,
the intended interpretation for <P⊃Q> is "if P, then Q equivalently, "P implies Q".

                        The Converse of the Fantasy Rule
Now Lewis Carroll's Dialogue was all about "if-then" statements. In particular, Achilles
had a lot of trouble in persuading the Tortoise to accept the second clause of an "if-then"
statement, even when the "if-then" state itself was accepted, as well as its first clause. The
next rule allows y infer the second "clause" of a'⊃'-string, provided that the `⊃'-string it a
theorem, and that its first "clause" is also a theorem.




The Propositional Calculus                                                                193

RULE OF DETACHMENT: If x and < x⊃y> are both theorems, then y is a theorem.

Incidentally, this rule is often called "Modus Ponens", and the fantasy rule is often called
the "Deduction Theorem".

                   The Intended Interpretation of the Symbols
We might as well let the cat out of the bag at this point, and reveal the "meanings" of the
rest of the symbols of our new system. In case it is not yet apparent, the symbol `A' is
meant to be acting isomorphically to the normal, everyday word `and'. The symbol '-'
represents the word 'not'-it is a formal sort of negation. The angle brackets '<' and `>' are
groupers-their function being very similar to that of parentheses in ordinary algebra. The
main difference is that in algebra, you have the freedom to insert parentheses or to leave
them out, according to taste and style, whereas in a formal system, such anarchic freedom
is not tolerated. The symbol '∨' represents the word `or' ('vel' is a Latin word for `or'). The
`or' that is meant is the so-called inclusive `or', which means that the interpretation of
<x∨y> is "either x or y-or both".

        The only symbols we have not interpreted are the atoms. An atom has no single
interpretation-it may be interpreted by any sentence of English (it must continue to be
interpreted by the same sentence if it occurs multiply within a string or derivation). Thus,
for example, the well-formed string <P∧~P> could be interpreted by the compound
sentence

This mind is Buddha, and this mind is not Buddha.

Now let us look at each of the theorems so far derived, and interpret them. The first one
was <P⊃~~P>. If we keep the same interpretation for P, we have the following
interpretation:

If this mind is Buddha,
         then it is not the case that this mind is not Buddha.

Note how I rendered the double negation. It is awkward to repeat a negation in any
natural language, so one gets around it by using two different ways of expressing
negation. The second theorem we derived was <<P∧Q>⊃<Q∧P>>. If we let Q be
interpreted by the sentence "This flax weighs three pounds", then our theorem reads as
follows:

If this mind is Buddha and this flax weighs three pounds,
         then this flax weighs three pounds and this mind is Buddha.

The third theorem was <P⊃<Q⊃<P∧Q>>>. This one goes into the following nested "if-
then" sentence:



The Propositional Calculus                                                                194

If this mind is Buddha,
         then, if this flax weighs three pounds,
                  then this mind is Buddha and this flax weighs three pounds.

        You probably have noticed that each theorem, when interpreted, something
absolutely trivial and self-evident. (Sometimes they are so s evident that they sound
vacuous and-paradoxically enough-confusing or even wrong!) This may not be very
impressive, but just remember there are plenty of falsities out there which could have
been produced they weren't. This system-the Propositional Calculus-steps neatly ft truth
to truth, carefully avoiding all falsities, just as a person who is concerned with staying dry
will step carefully from one stepping-stone creek to the next, following the layout of
stepping-stones no matter I twisted and tricky it might be. What is impressive is that-in
the Propositional Calculus-the whole thing is done purely typographically. There is
nobody down "in there", thinking about the meaning of the strings. It i! done
mechanically, thoughtlessly, rigidly, even stupidly.

                          Rounding Out the List of Rules
We have not yet stated all the rules of the Propositional Calculus. The complete set of
rules is listed below, including the three new ones.

JOINING RULE: If x and y are theorems, then < x∧y> is a theorem.

SEPARATION RULE: If < x∧y> is a theorem, then both x and y are theorems.

DOUBLE-TILDE RULE: The string '~~' can be deleted from any theorem can also be
  inserted into any theorem, provided that the result string is itself well-formed.

FANTASY RULE: If y can be derived when x is assumed to be a theorem then < x⊃y> is
  a theorem.

CARRY-OVER RULE: Inside a fantasy, any theorem from the "reality" c level higher
  can be brought in and used.

RULE OF DETACHMENT: If x and < x⊃y> are both theorems, then y is a theorem.

CONTRAPOSITIVE RULE: <x⊃y> and <~y⊃~x> are interchangeable

DE MORGAN'S RULE: <~x∧~y> and ~< x∨y> are interchangeable.

SWITCHEROO RULE: <x∨y> and <~x⊃y> are interchangeable.

(The Switcheroo rule is named after Q. q. Switcheroo, an Albanian railroad engineer who
worked in logic on the siding.) By "interchangeable" in foregoing rules, the following is
meant: If an expression of one form occurs as either a theorem or part of a theorem, the
other form may be


The Propositional Calculus                                                                195

substituted, and the resulting string will also be a theorem. It must be kept in mind that
the symbols ‘x’ and ‘y’ always stand for well-formed strings of the system.

                                 Justifying the Rules
Before we see these rules used inside derivations, let us look at some very short
justifications for them. You can probably justify them to yourself better than my
examples – which is why I only give a couple.
        The contrapositive rule expresses explicitly a way of turning around conditional
statements which we carry out unconsciously. For instance, the “Zentence”

       If you are studying it, then you are far from the Way

Means the same thing as

       If you are close to the Way, then you are not studying it.

        De Morgan’s rule can be illustrated by our familiar sentence “The flag is not
moving and the wind is not moving”. If P symbolizes “the flag is not moving”, and Q
symbolizes “the wind is moving”, then the compound sentence is symbolized by
<~P∧~Q>, which, according to Morgan’s law, is interchangeable with ~<P∨Q>. whose
interpretation would be “It is not true that either the flag or the wind is moving”. And no
one could deny that it is a Zensible conclusion to draw.
        For the Switrcheroo rule, consider the sentence “Either a cloud is hanging over
the mountain, or the moonlight is penetrating the waves of the lake,” which might be
spoken, I suppose, by a wistful Zen master remembering a familiar lake which he can
visualize mentally but cannot see. Now hang on to your seat, for the Swircheroo rule tells
us that this is interchangeable with the thought “If a cloud is not hanging over the
mountain, then the moonlight is penetrating the waves of the lake.” This may not be
enlightenment, but it is the best the Propositional Calculus has to offer.

                          Playing around with the system
Now, let us apply these rules to a previous theorem, ands see what we get: For instance,
take the theorem <P⊃~~P>:

<P⊃~~P>:                      old theorem
<~~~P⊃~P>:                    contrapositive
<~P⊃~P>                       double-tilde
<P∨~P>                        switcheroo

This new theorem, when interpreted, says:




The Propositional Calculus                                                             196

       Either this mind is Buddha, or this mind is not Buddha

Once again, the interpreted theorem, though perhaps less than mind boggling, is at least
true.

                                 Semi-Interpretations
It is natural, when one reads theorems of the Propositional Calculus out loud, to interpret
everything but the atoms. I call this semi-interpreting. For example, the semi-
interpretation of <P∨~P>:: would be

                                        P or not P.

Despite the fact that P is not a sentence, the above semisentence still sounds true, because
you can very easily imagine sticking any sentence in for P – and the form of the semi-
interpreted theorem assures you that however you make your choice, the resulting
sentence will be true. And that is the key idea of the Propositional Calculus: it produces
theorems which, when semi-interpreted, are seen to be “universally true semisaentences”,
by which is meant that no matter how you complete the interpretation, the final result will
be a true statement.

                                       Ganto’s Ax
Now we can do a more advanced exercise, based on a Zen koan called “Ganto’s Ax”.
Here is how it began.

   One day Tokusan told his student Ganto, “I have two monks who have been here
   for many years. Go and examine them.” Ganto picked up an ax and went to the hut
   where the two monks were meditating. He raised the ax, saying “If you say a word,
   I will cut off your heads; and if you do not say a word, I will also cut off your
   heads.”1

If you say a word I will cut off this koan, and if you do not say a word, I will also cut off
this koan – because I want you to translate some of it into our notation. Let us symbolize
“you say a word” by P and “I will cut off your heads” by Q. Then Ganto’s ax threat is
symbolized by the string <<P⊃Q>∧<~`P⊃Q>>. What if this ax threat were an axiom?
Here is a fantasy to answer that question.

(1) [                                         push
(2)   <<P⊃Q>∧<~`P⊃Q>>.                        Ganto’s axiom
(3)   <P⊃Q>                                   separation
(4)   <~Q⊃~P>.                                contrapositive
(5)   <~P⊃Q>                                  separation
(6)   <~Q⊃~~P>.                               contrapositive
(7)    ]                                      push again
(8)      ~Q                                   premise

The Propositional Calculus                                                               197

(9)      <~Q⊃~P>.                     carry-over of line 4
(10)      ~P                          detachment
(11)     <~Q⊃~~P>.                    carry-over of line 6
(12)      ~~P                         detachment (lines 8 and 11)
(13)     <~P∧~~P>                     joining
(14)      <~P∨~~P>                    De Morgan
(15)   ]                              pop once
(16)   <~Q⊃~<P∨~P>>.                  fantasy rule
(17)   <~P∨~P>⊃Q>.                    contrapositive
(18) [                                push
(19)   . ~P                           premise (also outcome)
(20) ]                                pop
(21)   <~P⊃~P>.                       fantasy rule
(22)   <P∨~P>.                        switcheroo
(23)   Q                              detachment (lines 22 and 17)
(24) ]                                pop out

The power of the Propositional Calculus is shown in this example. Why, in but two dozen
steps, we have deduced Q: that the heads will be cut off! (Ominously, the rule last
invoked was "detachment" ...) It might seem superfluous to continue the koan now, since
we know what must ensue ... However, I shall drop my resolve to cut the koan off; it is a
true Zen koan, after all. The rest of the incident is here related:

   Both monks continued their meditation as if he had not spoken. Ganto dropped the
   ax and said, "You are true Zen students." He returned to Tokusan and related the
   incident. "I see your side well," Tokusan agreed, "but tell me, how is their side?"
   "Tõzan may admit them," replied Ganto, "but they should not be admitted under
   Tokusan."2

Do you see my side well? How is the Zen side?

                 Is There a Decision Procedure for Theorems?
The Propositional Calculus gives us a set of rules for producing statements which would
be true in all conceivable worlds. That is why all of its theorems sound so simple-minded;
it seems that they have absolutely no content! Looked at this way, the Propositional
Calculus might seem to be a waste of time, since what it tells us is absolutely trivial. On
the other hand, it does it by specifying the form of statements that are universally true,
and this throws a new kind of light onto the core truths of the universe: they are not only
fundamental, but also regular: they can be produced by one set of typographical rules. To
put it another way, they are all "cut from the same cloth". You might consider whether
the same could be said about Zen koans: could they all be produced by one set of
typographical rules?
        It is quite relevant here to bring up the question of a decision procedure. That is,
does there exist any mechanical method to tell nontheorems from theorems? If so, that
would tell us that the set of theorems of the

The Propositional Calculus                                                              198

Propositional Calculus is not only r.e., but also recursive. It turns out that there is an
interesting decision procedure-the method of truth u would take us a bit afield to present
it here; you can find it in almost any standard book on logic. And what about Zen koans?
Could there conceivably be a mechanical decision procedure which distinguishes genuine
Zen koans from other things?

                     Do We Know the System Is Consistent?
Up till now, we have only presumed that all theorems, when interpreted as indicated, are
true statements. But do we know that that is the case' we prove it to be? This is just
another way of asking whether the intended interpretations ('and' for `∧', etc.) merit being
called the "passive meanings” of the symbols. One can look at this issue from two very
different points of view, which might be called the "prudent" and "imprudent" points I
will now present those two sides as I see them, personifying their as "Prudence" and
"Imprudence".

Prudence: We will only KNOW that all theorems come out true un intended
   interpretation if we manage to PROVE it. That is the c: thoughtful way to proceed.
Imprudence: On the contrary. It is OBVIOUS that all theorems will come out true. If you
   doubt me, look again at the rules of the system. You will find that each rule makes a
   symbol act exactly as the word it represents ought to be used. For instance, the joining
   rule makes the symbol ‘∧’ act as `and' ought to act; the rule of detachment makes `⊃'
   act as it ought to, if it is to stand for 'implies', or 'if-then'; and so on. Unless you are
   like the Tortoise, you will recognize in each rule a codification of a pattern you use in
   your own thought patterns. So if you trust your own thought patterns, then you HAVE
   to believe that all theorems come out true! That's the way I see it. I don't need any
   further proof. If you think that some theorem comes out false, then presumably you
   think that some rule must be wrong. Show me which one.
Prudence: I'm not sure that there is any faulty rule, so I can't point one out to you. Still, I
   can imagine the following kind of scenario. You, following the rules, come up with a
   theorem -- say x. Meanwhile I, also following the rules, come up with another
   theorem-it happens to be ~x.           Can't you force yourself to conceive of that?
Imprudence: All right; let's suppose it happened. Why would it bother you? Or let me put
   it another way. Suppose that in playing with the MIU-system, I came up with a
   theorem x, and you came up with xU Can you force yourself to conceive of that?
Prudence: Of course-in fact both MI and MIU are theorems.
Imprudence: Doesn't that bother you?
Prudence: Of course not. Your example is ridiculous, because MI and MIU are not
   CONTRADICTORY, whereas two strings x and ~x in the Propositional Calculus
   ARE contradictory.




The Propositional Calculus                                                                 199

Imprudence: Well, yes -- provided you wish to interpret `~' as `not'. But what would lead
   you to think that '~' should be interpreted as `not'?
Prudence: The rules themselves. When you look at them, you realize that the only
   conceivable interpretation for '~' is 'not'-and likewise, the only conceivable
   interpretation for `∧' is `and', etc.
Imprudence: In other words, you are convinced that the rules capture the meanings of
   those words?
Prudence: Precisely.
Imprudence: And yet you are still willing to entertain the thought that both x and ~x
   could be theorems? Why not also entertain the notion that hedgehogs are frogs, or that
   1 equals 2, or that the moon is made of green cheese? I for one am not prepared even
   to consider whether such basic ingredients of my thought processes are wrong --
   because if I entertained that notion, then I would also have to consider whether my
   modes of analyzing the entire question are also wrong, and I would wind up in a total
   tangle.
Prudence: Your arguments are forceful ... Yet I would still like to see a PROOF that all
   theorems come out true, or that x and ~x can never both be theorems.
Imprudence: You want a proof. I guess that means that you want to be more convinced
   that the Propositional Calculus is consistent than you are convinced of your own
   sanity. Any proof I could think of would involve mental operations of a greater
   complexity than anything in the Propositional Calculus itself. So what would it prove?
   Your desire for a proof of consistency of the Propositional Calculus makes me think
   of someone who is learning English and insists on being given a dictionary which
   definers all the simple words in terms of complicated ones...

                             The Carroll Dialogue Again
This little debate shows the difficulty of trying to use logic and reasoning to defend
themselves. At some point, you reach rock bottom, and there is no defense except loudly
shouting, "I know I'm right!" Once again, we are up against the issue which Lewis
Carroll so sharply set forth in his Dialogue: you can't go on defending your patterns of
reasoning forever. There comes a point where faith takes over.
        A system of reasoning can be compared to an egg. An egg has a shell which
protects its insides. If you want to ship an egg somewhere, though, you don't rely on the
shell. You pack the egg in some sort of container, chosen according to how rough you
expect the egg's voyage to be. To be extra careful, you may put the egg inside several
nested boxes. However, no matter how many layers of boxes you pack your egg in, you
can imagine some cataclysm which could break the egg. But that doesn't mean that you'll
never risk transporting your egg. Similarly, one can never give an ultimate, absolute
proof that a proof in some system is correct. Of course,




The Propositional Calculus                                                           200

one can give a proof of a proof, or a proof of a proof of a proof – but the validity of the
outermost system always remains an unproven assumption, accepted on faith. One can
always imagine that some unsuspected subtlety will invalidate every single level of proof
down to the bottom, and tI "proven" result will be seen not to be correct after all. But that
doesn’t mean that mathematicians and logicians are constantly worrying that the whole
edifice of mathematics might be wrong. On the other hand, unorthodox proofs are
proposed, or extremely lengthy proofs, or proofs generated by computers, then people do
stop to think a bit about what they really mean by that quasi-sacred word "proven".
        An excellent exercise for you at this point would be to go back Carroll Dialogue,
and code the various stages of the debate into our notation -- beginning with the original
bone of contention:

Achilles: If you have <<A∧B>⊃Z>, and you also have <A∧B>, then surely you have Z.
Tortoise: Oh! You mean: <<<<A∧B>⊃Z>∧<A∧B>>⊃Z>, : don't you?

(Hint: Whatever Achilles considers a rule of inference, the Tortoise immediately flattens
into a mere string of the system. If you use or letters A, B, and Z, you will get a recursive
pattern of longer and strings.)

                            Shortcuts and Derived Rules
When carrying out derivations in the Propositional Calculus, one quickly invents various
types of shortcut, which are not strictly part of the system For instance, if the string
<Q∨~Q> were needed at some point, and <P∨~P> had been derived earlier, many people
would proceed as if <Q∨~Q> had been derived, since they know that its derivation is an
exact parallel to that of <P∨~P>. The derived theorem is treated as a "theorem schema" --
a mold for other theorems. This turns out to be a perfect valid procedure, in that it always
leads you to new theorems, but it is not a rule of the Propositional Calculus as we
presented it. It is, rather, a derived rule, It is part of the knowledge which we have about
the system. That this rule keeps you within the space of theorems needs proof, of course-
but such a proof is not like a derivation inside the system. It is a proof in the ordinary,
intuitive sense -- a chain of reasoning carried out in the I-mode. The theory about the
Propositional Calculus is a "metatheory", and results in it can be called "metatheorems" -
Theorems about theorems. (Incidentally, note the peculiar capitalization in the phrase
"Theorems about theorems". It is a consequence of our convention: metatheorems are
Theorems (proven results) concerning theorems (derivable strings).)
        In the Propositional Calculus, one could discover many metatheorems, or derived
rules of inference. For instance, there is a De Morgan's Rule:




The Propositional Calculus                                                               201

                        <~x∨~y> and ~<x∧y> are interchangeable.

If this were a rule of the system, it could speed up many derivations considerably. But if
we prove that it is correct, isn't that good enough? Can't we use it just like a rule of
inference, from then on?
        There is no reason to doubt the correctness of this particular derived rule. But
once you start admitting derived rules as part of your procedure in the Propositional
Calculus, you have lost the formality of the system, since derived rules are derived
informally-outside the system. Now formal systems were proposed as a way to exhibit
every step of a proof explicitly, within one single, rigid framework, so that any
mathematician could check another's work mechanically. But if you are willing to step
outside of that framework at the drop of a hat, you might as well never have created it at
all. Therefore, there is a drawback to using such shortcuts.

                             Formalizing Higher Levels
On the other hand, there is an alternative way out. Why not formalize the metatheory,
too? That way, derived rules (metatheorems) would be theorems of a larger formal
system, and it would be legitimate to look for shortcuts and derive them as theorems-that
is, theorems of the formalized metatheory-which could then be used to speed up the
derivations of theorems of the Propositional Calculus. This is an interesting idea, but as
soon as it is suggested, one jumps ahead to think of metametatheories, and so on. It is
clear that no matter how many levels you formalize, someone will eventually want to
make shortcuts in the top level.
         It might even be suggested that a theory of reasoning could be identical to its own
metatheory, if it were worked out carefully. Then, it might seem, all levels would
collapse into one, and thinking about the system would be just one way of working in the
system! But it is not that easy. Even if a system can "think about itself", it still is not
outside itself. You, outside the system, perceive it differently from the way it perceives
itself. So there still is a metatheory-a view from outside-even for a theory which can
"think about itself" inside itself. We will find that there are theories which can "think
about themselves". In fact, we will soon see a system in which this happens completely
accidentally, without our even intending it! And we will see what kinds of effects this
produces. But for our study of the Propositional Calculus, we will stick with the simplest
ideas-no mixing of levels.
         Fallacies can result if you fail to distinguish carefully between working in the
system (the M-mode) and thinking about the system (the I-mode). For example, it might
seem perfectly reasonable to assume that, since <P∨~P> (whose semi-interpretation is
"either P or not P") is a theorem, either P or ~P must be a theorem. But this is dead
wrong: neither one of the latter pair is a theorem. In general, it is a dangerous practice to
assume that symbols can be slipped back and forth between different levels-here, the
language of the formal system and its metalanguage (English).




The Propositional Calculus                                                               202

        Reflections on the Strengths and Weaknesses of the System
You have now seen one example of a system with a purpose-to re part of the architecture
of logical thought. The concepts which this handles are very few in number, and they are
very simple, precise co But the simplicity and precision of the Propositional Calculus are
the kinds of features which make it appealing to mathematicians. There are two reasons
for this. (1) It can be studied for its own properties, ex geometry studies simple, rigid
shapes. Variants can be made on it, employing different symbols, rules of inference,
axioms or axiom schemata on. (Incidentally, the version of the Propositional Calculus
here pr is related to one invented by G. Gentzen in the early 1930's. The other versions in
which only one rule of inference is used-detachment usually-and in which there are
several axioms, or axiom schemata study of ways to carry out propositional reasoning in
elegant formal systems is an appealing branch of pure mathematics. (2) The Propositional
Calculus can easily be extended to include other fundamental aspects of reasoning. Some
of this will be shown in the next Chapter, where the Propositional Calculus is
incorporated lock, stock and barrel into a much larger and deeper system in which
sophisticated number-theoretical reasoning can be done.

                                Proofs vs. Derivations
The Propositional Calculus is very much like reasoning in some w one should not equate
its rules with the rules of human thought. A proof is something informal, or in other
words a product of normal thought written in a human language, for human consumption.
All sorts of complex features of thought may be used in proofs, and, though they may
“feel right", one may wonder if they can be defended logically. That is really what
formalization is for. A derivation is an artificial counterpart of and its purpose is to reach
the same goal but via a logical structure whose methods are not only all explicit, but also
very simple.
        If -- and this is usually the case -it happens that a formal derivation is extremely
lengthy compared with the corresponding "natural" proof that is just too bad. It is the
price one pays for making each step so simple. What often happens is that a derivation
and a proof are "simple" in complementary senses of the word. The proof is simple in
that each step sounds right", even though one may not know just why; the derivation is
simple in that each of its myriad steps is considered so trivial that it is beyond reproach,
and since the whole derivation consists just of such trivial steps it is supposedly error-
free. Each type of simplicity, however, brings along a characteristic type of complexity.
In the case of proofs, it is the complexity of the underlying system on which they rest --
namely, human language -- and in the case of derivations, it is their astronomical size,
which makes them almost impossible to grasp.
        Thus, the Propositional Calculus should be thought of as part of a




The Propositional Calculus                                                                203

general method for synthesizing artificial proof-like structures. It does not, however, have
much flexibility or generality. It is intended only for use in connection with mathematical
concepts-which are themselves quite rigid. As a rather interesting example of this, let us
make a derivation in which a very peculiar string is taken as a premise in a fantasy:
<P∧~P>. At least its semi-interpretation is peculiar. The Propositional Calculus,
however, does not think about semi-interpretations; it just manipulates strings
typographically-and typographically, there is really nothing peculiar about this string.
Here is a fantasy with this string as its premise:

       (1)     [                      push
       (2)         <P∧~P>             premise
        (3)        P                  separation
       (4)         ~P                 separation
       (5)         [                  push
       (6)            ~Q              premise
       (7)           P                carry-over line 3
       (8)           ~~P              double-tilde
       (9)          ]                 pop
       (10)         <~Q⊃~~P>          fantasy
       (11)         <~P⊃Q>            contrapositive
       (12)        Q                  detachment (Lines 4,11)
       (13)    ]                      pop
       (14)    <<P∧~P >⊃Q>            fantasy

Now this theorem has a very strange semi-interpretation:

       P and not P together imply Q

Since Q is interpretable by any statement, we can loosely take the theorem to say that
"From a contradiction, anything follows"! Thus, in systems based on the Propositional
Calculus, contradictions cannot be contained; they infect the whole system like an
instantaneous global cancer.

                         The Handling of Contradictions
This does not sound much like human thought. If you found a contradiction in your own
thoughts, it's very unlikely that your whole mentality would break down. Instead, you
would probably begin to question the beliefs or modes of reasoning which you felt had
led to the contradictory thoughts. In other words, to the extent you could, you would step
out of the systems inside you which you felt were responsible for the contradiction, and
try to repair them. One of the least likely things for you to do would be to throw up your
arms and cry, "Well, I guess that shows that I believe everything now!" As a joke, yes-
but not seriously.
         Indeed, contradiction is a major source of clarification and progress in all domains
of life-and mathematics is no exception. When in times past, a


The Propositional Calculus                                                               204

contradiction in mathematics was found, mathematicians would immediately seek to
pinpoint the system responsible for it, to jump out of it, to reason about it, and to amend
it. Rather than weakening mathematics, the discovery and repair of a contradiction would
strengthen it. This might take time and a number of false starts, but in the end it would
yield fruit. For instance, in the Middle Ages, the value of the infinite series

       1 – 1 + 1 – 1 + 1 -. ..

was hotly disputed. It was "proven" to equal 0, 1, ½, and perhaps other values. Out of
such controversial findings came a fuller, deeper about infinite series.
        A more relevant example is the contradiction right now confronting us-namely the
discrepancy between the way we really think, and t the Propositional Calculus imitates
us. This has been a source of discomfort for many logicians, and much creative effort has
gone into trying to patch up the Propositional Calculus so that it would not act so stupidly
and inflexibly. One attempt, put forth in the book Entailment by A. R. Anderson and N.
Belnap,3 involves "relevant implication", which tries to make the symbol for "if-then"
reflect genuine causality, or at least connect meanings. Consider the following theorems
of the Propositional Calculus

                                         <P⊃<Q⊃P>>
                                         <P⊃<Q∨~P>>
                                         <<P∧~P>⊃Q>
                                       <<P⊃Q>∨<Q⊃P>>


They, and many others like them, all show that there need be no relationship at all
between the first and second clauses of an if-then statement for it to be provable within
the Propositional Calculus. In protest, "relevant implication" puts certain restrictions on
the contexts in which the rules of inference can be applied. Intuitively, it says that
"something can only be derived from something else if they have to do with each other”.
For example, line 10 in the derivation given above would not be allowed in such a
system, and that would block the derivation of the <<P∧~P >⊃Q>
        More radical attempts abandon completely the quest for completeness or
consistency, and try to mimic human reasoning with all its inconsistencies. Such research
no longer has as its goal to provide a solid underpinning for mathematics, but purely to
study human thought processes.
        Despite its quirks, the Propositional Calculus has some feat recommend itself. If
one embeds it into a larger system (as we will do next Chapter), and if one is sure that the
larger system contains no contradictions (and we will be), then the Propositional Calculus
does all that one could hope: it provides valid propositional inferences -- all that can be
made. So if ever an incompleteness or an inconsistency is uncovered, can be sure that it
will be the fault of the larger system, and not of its subsystem which is the Propositional
Calculus.




The Propositional Calculus                                                              205

             FIGURE 42. “Crab Canon”, by M. C. Escher (~1965)




Crab Canon                                                      206

